% ================================================================
%  《收获与播种》统一封面 —— TikZ 原生矢量绘制(不依赖任何外部图片)
%  母题: 暮色中的麦田 —— 犁沟向地平线的灭点收束, 前景麦穗剪影,
%        一位播种者把一把种子撒向天空, 种子沿弧线落入田里。
%  用法(在 \begin{document} 之前, 依次):
%      % ================================================================
%  《收获与播种》中文译本 —— 各卷共用导言区
%  由 vol-prelude / vol-part1 / vol-part2 / vol-part3 / vol-part4
%  以及全书合订 main-all 共同 \input。基类 mathtranslation.cls v3.2 不改。
%  --------------------------------------------------------------
%  样式均在此改写, 以贴近原书:
%    · 章标题仿原书 "5. Travail et découverte"(块状左对齐黑体)
%    · 节/小节: 黑体左对齐, 编号后带句点
%    · 脚注: 人工编号照录原书——前奏与第一部分页内重新编号,
%      第二至四部各自连续编号并带(*)角标; 第一部分另有 \fnoteP
%      处理 "(n)" 式连续注号
%    · 边码: 原打字稿页码 \origpage{N}(前奏作 A1、P1、L1、罗马数字等)
%    · ⋄ 记号: 段首补写标记 \pmark 与上标形式 \dmark
% ================================================================

% 题名含空格, 不能走类选项(会被 kvoptions 吃掉), 在此赋值
\renewcommand{\BookTitleEN}{Récoltes et Semailles}
\renewcommand{\OriginalAuthor}{Alexandre Grothendieck}
\renewcommand{\TranslationDate}{\today}

% --------------------------------------------------------------
%  罕用字后备(如"祐"): Fandol 不载时退到 SimSun
% --------------------------------------------------------------
\setCJKfallbackfamilyfont{\CJKrmdefault}{{SimSun}}

% --------------------------------------------------------------
%  hyperref 元信息(各卷在 % ================================================================
%  《收获与播种》中文译本 —— 各卷共用导言区
%  由 vol-prelude / vol-part1 / vol-part2 / vol-part3 / vol-part4
%  以及全书合订 main-all 共同 \input。基类 mathtranslation.cls v3.2 不改。
%  --------------------------------------------------------------
%  样式均在此改写, 以贴近原书:
%    · 章标题仿原书 "5. Travail et découverte"(块状左对齐黑体)
%    · 节/小节: 黑体左对齐, 编号后带句点
%    · 脚注: 人工编号照录原书——前奏与第一部分页内重新编号,
%      第二至四部各自连续编号并带(*)角标; 第一部分另有 \fnoteP
%      处理 "(n)" 式连续注号
%    · 边码: 原打字稿页码 \origpage{N}(前奏作 A1、P1、L1、罗马数字等)
%    · ⋄ 记号: 段首补写标记 \pmark 与上标形式 \dmark
% ================================================================

% 题名含空格, 不能走类选项(会被 kvoptions 吃掉), 在此赋值
\renewcommand{\BookTitleEN}{Récoltes et Semailles}
\renewcommand{\OriginalAuthor}{Alexandre Grothendieck}
\renewcommand{\TranslationDate}{\today}

% --------------------------------------------------------------
%  罕用字后备(如"祐"): Fandol 不载时退到 SimSun
% --------------------------------------------------------------
\setCJKfallbackfamilyfont{\CJKrmdefault}{{SimSun}}

% --------------------------------------------------------------
%  hyperref 元信息(各卷在 % ================================================================
%  《收获与播种》中文译本 —— 各卷共用导言区
%  由 vol-prelude / vol-part1 / vol-part2 / vol-part3 / vol-part4
%  以及全书合订 main-all 共同 \input。基类 mathtranslation.cls v3.2 不改。
%  --------------------------------------------------------------
%  样式均在此改写, 以贴近原书:
%    · 章标题仿原书 "5. Travail et découverte"(块状左对齐黑体)
%    · 节/小节: 黑体左对齐, 编号后带句点
%    · 脚注: 人工编号照录原书——前奏与第一部分页内重新编号,
%      第二至四部各自连续编号并带(*)角标; 第一部分另有 \fnoteP
%      处理 "(n)" 式连续注号
%    · 边码: 原打字稿页码 \origpage{N}(前奏作 A1、P1、L1、罗马数字等)
%    · ⋄ 记号: 段首补写标记 \pmark 与上标形式 \dmark
% ================================================================

% 题名含空格, 不能走类选项(会被 kvoptions 吃掉), 在此赋值
\renewcommand{\BookTitleEN}{Récoltes et Semailles}
\renewcommand{\OriginalAuthor}{Alexandre Grothendieck}
\renewcommand{\TranslationDate}{\today}

% --------------------------------------------------------------
%  罕用字后备(如"祐"): Fandol 不载时退到 SimSun
% --------------------------------------------------------------
\setCJKfallbackfamilyfont{\CJKrmdefault}{{SimSun}}

% --------------------------------------------------------------
%  hyperref 元信息(各卷在 \input{common} 之后可用 \hypersetup 覆写)
% --------------------------------------------------------------
\hypersetup{
  pdftitle={收获与播种},
  pdfauthor={Alexandre Grothendieck / 数译 译},
  pdfsubject={Récoltes et Semailles — Réflexions et témoignage sur un passé de mathématicien}
}

% --------------------------------------------------------------
%  标题样式: 仿原书 —— 黑体左对齐, 编号后带句点
% --------------------------------------------------------------
\titleformat{\section}[block]
  {\normalfont\Large\headfont}{\thesection.}{0.5em}{}
\titlespacing*{\section}{0pt}{3.2ex plus 1.2ex minus .2ex}{1.6ex plus .2ex}
\titleformat{\subsection}[block]
  {\normalfont\large\headfont}{\thesubsection.}{0.5em}{}
\titlespacing*{\subsection}{0pt}{2.6ex plus 1ex minus .2ex}{1.2ex plus .2ex}
\titleformat{\subsubsection}[block]
  {\normalfont\normalsize\headfont}{\thesubsubsection.}{0.5em}{}
\titlespacing*{\subsubsection}{0pt}{2.2ex plus .8ex minus .2ex}{1ex plus .2ex}
\titleformat{name=\subsubsection,numberless}[block]
  {\normalfont\normalsize\headfont}{}{0em}{}
\setcounter{secnumdepth}{3}
\setcounter{tocdepth}{3}
% 为含长外文书名与行内公式的段落提供必要的断行余量
\setlength{\emergencystretch}{3em}

% 章标题样式: 原书章标题与节同级排印("13. A) HERITAGE..."),
% 故章标题取块状左对齐黑体, 字号介于节与部之间
\titleformat{\chapter}[block]
  {\normalfont\LARGE\headfont}{\thechapter.}{0.5em}{}
\titlespacing*{\chapter}{0pt}{3.6ex plus 1.4ex minus .2ex}{2ex plus .2ex}
% 目录中章号照原书用阿拉伯数字("13. A) ..." 而非 "第十三章")
\renewcommand{\CTEXthechapter}{\thechapter\hspace{.3em}}

% --------------------------------------------------------------
%  目录观感: 条目与引导点之间、引导点与页码之间各留 0.5em
%  内核 \@dottedtocline 让 \leaders 从标题末字起排, 第一个点几乎
%  贴着末字, 读起来像多标了一个句点; 页码前也没有余量。两处各补
%  0.5em。
% --------------------------------------------------------------
\makeatletter
\patchcmd{\@dottedtocline}{\nobreak\leaders}{\nobreak\hspace{0.5em}\leaders}%
  {\typeout{MT: TOC 引导点前已留 0.5em}}%
  {\typeout{MT-WARN: TOC 引导点前间距补丁未生效}}
\patchcmd{\@dottedtocline}{\hfill\nobreak\hb@xt@\@pnumwidth}%
  {\hfill\nobreak\hspace{0.5em}\hb@xt@\@pnumwidth}%
  {\typeout{MT: TOC 页码前已留 0.5em}}%
  {\typeout{MT-WARN: TOC 页码前间距补丁未生效}}
\makeatother

% --------------------------------------------------------------
%  正文半角标点后的间距
%  本书体例用半角 , ; : ? ! 与 ASCII 括号, 由 xeCJK 按 HalfRight
%  (右半留白) 处理。放宽到 0.45em, 与全角标点右半留白大致齐平。
% --------------------------------------------------------------
\xeCJKsetup{CJKecglue={\hskip 0.45em plus 0.06em minus 0.09em}}

% --------------------------------------------------------------
%  页眉页脚: 偶数页眉右对齐斜体节题; 页码页脚外侧
% --------------------------------------------------------------
\fancyhf{}
\fancyhead[RE]{\small\itshape\rightmark}
\fancyfoot[LE,RO]{\small\thepage}
\renewcommand{\headrulewidth}{0pt}
\renewcommand{\sectionmark}[1]{\markright{\thesection.\ #1}}
\fancypagestyle{plain}{%
  \fancyhf{}%
  \fancyfoot[LE,RO]{\small \thepage}%
  \renewcommand{\headrulewidth}{0pt}%
}

% --------------------------------------------------------------
%  原书版式记号宏
% --------------------------------------------------------------
% 边码: 原打字稿页码(G. 手稿系统里的自引页码), 排在页边
\usepackage{marginnote}
\newcommand{\origpage}[1]{\marginnote{\normalfont\scriptsize p.~#1}}

% ⋄ 记号: 段首/行内正常大小; 上标形式为作者后补文字的位置标记
\DeclareRobustCommand{\pmark}{\ensuremath{\diamond}\hspace{0.15em}}
\DeclareRobustCommand{\dmark}{\textsuperscript{\ensuremath{\diamond}}}

% 脚注人工编号: \fnote{795}{(*)}{文本} —— 编号与(*)角标均照录原书
\newcommand{\fnote}[3]{%
  \begingroup
  \renewcommand{\thefootnote}{#1#2}%
  \footnote{#3}%
  \endgroup
}

% 第一部分的脚注: 页内序号作上标(#1), 注文行首印 "序号(连续注号)",
% 正文在序号之后另行括注同一连续注号 —— \fnoteP{1}{(3)}{注文}
% 原书有三条注无连续注号而改标日期(ch07 注3 = 30 Septembre,
% ch09 注11/12 = 8 août), 第二参数留空, 日期写进注文之首。
\newcommand{\fnoteP}[3]{%
  \begingroup
  \renewcommand{\thefootnote}{#1}\footnotemark
  \renewcommand{\thefootnote}{#1#2}\footnotetext{#3}%
  \endgroup
}

% 正文呼号与注文行首标号不同(原书照录)
\newcommand{\fnotevariant}[4]{%
  \begingroup
  \renewcommand{\thefootnote}{#1#2}\footnotemark
  \renewcommand{\thefootnote}{#1#3}\footnotetext{#4}%
  \endgroup
}

% 大注块标题: \notehead{171³}{斜体评论} —— 正文随后另起
\newcommand{\notehead}[2]{%
  \par\bigskip
  \noindent{\headfont 注~#1}\ifblank{#2}{}{\hspace{0.5em}{\itshape #2}}%
  \par\nobreak\smallskip\nobreak
}

% ∗ ∗ ∗ 分隔
\newcommand{\asterisms}{%
  \par\vspace{1.2ex}%
  \noindent\hspace*{2em}$\ast$\hspace{1.2em}$\ast$\hspace{1.2em}$\ast$%
  \par\vspace{1.2ex}%
}

% 无编号子目的目录条目
\newcommand{\tocsubsub}[1]{%
  \phantomsection
  \addcontentsline{toc}{subsubsection}{#1}%
}

% 部扉: 仿原书部题页
\newcommand{\partdivider}[3]{%
  \clearpage
  \thispagestyle{empty}%
  \phantomsection
  \addcontentsline{toc}{part}{#1\quad #3}%
  \vspace*{6.5cm}
  \begin{center}
    {\fontsize{16pt}{26pt}\selectfont\headfont #1\par}
    \vspace{0.9cm}
    {\fontsize{22pt}{34pt}\selectfont\headfont #2\par}
  \end{center}
  \clearpage
}

% 献辞页: 仿原书部题后的题献, 独占一页且不入目录
\newcommand{\dedication}[1]{%
  \thispagestyle{empty}%
  \vspace*{4cm}%
  \begin{center}#1\end{center}%
  \clearpage
}

% --------------------------------------------------------------
%  各卷"译者说明"的共用体例段(各卷的范围段自写, 后接本段)
% --------------------------------------------------------------
\newcommand{\CommonNoteBody}{%
为便于与原书及学界通行的打字稿页码对勘, 本册保留:

\begin{itemize}
  \item 原书全部章节编号(章~1--18, 节如~5.2、18.2, 子目如~18.1.1)与原书小节标题;
  \item 原书脚注的编号与角标——
        \begin{itemize}
          \item 前奏(第~1--4 章)与第一部的注在各页内重新编号, 正文中的上标即该
                页序号; 第一部紧随其后的括注(如“(3)”)是该部分的连续注号, 与第~12
                章“第一部分的注”里的“注~3”一一对应, 注文行首照原书印作“1(3)”式。
                其中三条注原书未标连续注号而改标日期(ch7 注~3 = 30 Septembre、
                ch9 注~11/12 = 8 août), 本册照录其日期;
          \item 第二部第~13--16~章各自重新编号, 分别为~1--41、1--86、1--132、1--29;
                第三部第~17~章自~1~起至~12, 第~18~章自~1~重新起算并延续至~794;
                末段自~\textbf{795(*)}~续至~1053(*)。各段与原书一一对应;
        \end{itemize}
  \item 页边码“p.~1”“p.~A3”“p.~1066”式小注——这是作者引用自己打字稿原稿时
        所用的页码, 与本排印本页码不同。前奏四章的边码各有前缀: 第~1 章作
        \texttt{A1…A6}, 第~2 章作 \texttt{P1…P65}, 第~3 章作 \texttt{L1…L56},
        第~4 章作罗马数字 \texttt{I…XXII}, 自第一部起作纯数字 \texttt{1…171};
  \item 作者用以标记段落的~$\diamond$~记号, 以及标示后来补写文字的上标
        ~\textsuperscript{$\diamond$}~记号;
  \item 原书大小注(“注~44'”“注~46$^1$”“注~171³”等)的编号与层级。
\end{itemize}

第~12 章“第一部分的注”把第~5 至第~11 章各页脚注汇为一编, 并给每条注配上节题
(如“12.19. 耶稣与十二使徒”), 因此部分注文与正文脚注重复——这是原书的体例, 非
重出之误。这一部分属随笔式的自省文字, 同一句法语在正文脚注与章末汇编注中
各有措辞, 本册有意保留各自的译笔, 不强求统一。

除上述标记外, 译文力求贴近作者驳杂、跳跃而辛辣的文风; 数学名词的译法见各卷
术语表(如~\emph{topos}~不译, \emph{faisceaux pervers}~作“反常层”)。文中的日期
括注(如“(5~月~30~日)”“(1984年3月补写)”)为作者历次补写时的纪日, 一律照录。文中
另有一些对注号、页码的引用(如“n°~1695”“打字稿第~227~页”)沿袭作者打字稿的编号
系统, 与本排印本不尽一致, 亦无法在本书中落实, 悉照原文录出。

全套五卷(另有全书合订一册)的分工是:

\begin{itemize}
  \item \textbf{卷一·前奏}——主题的呈现: 或四个乐章的前奏(原书 p.23--122);
  \item \textbf{卷二·第一部}——愚妄与更新(原书 p.123--254);
  \item \textbf{卷三·第二部}——葬礼(I): 或中国皇帝的新衣(原书 p.255--412);
  \item \textbf{卷四·第三部}——葬礼(II): 或阴阳之钥(原书 p.413--632);
  \item \textbf{卷五·第四部}——葬礼(III): 或四种运算(原书 p.633--929)。
\end{itemize}
}
 之后可用 \hypersetup 覆写)
% --------------------------------------------------------------
\hypersetup{
  pdftitle={收获与播种},
  pdfauthor={Alexandre Grothendieck / 数译 译},
  pdfsubject={Récoltes et Semailles — Réflexions et témoignage sur un passé de mathématicien}
}

% --------------------------------------------------------------
%  标题样式: 仿原书 —— 黑体左对齐, 编号后带句点
% --------------------------------------------------------------
\titleformat{\section}[block]
  {\normalfont\Large\headfont}{\thesection.}{0.5em}{}
\titlespacing*{\section}{0pt}{3.2ex plus 1.2ex minus .2ex}{1.6ex plus .2ex}
\titleformat{\subsection}[block]
  {\normalfont\large\headfont}{\thesubsection.}{0.5em}{}
\titlespacing*{\subsection}{0pt}{2.6ex plus 1ex minus .2ex}{1.2ex plus .2ex}
\titleformat{\subsubsection}[block]
  {\normalfont\normalsize\headfont}{\thesubsubsection.}{0.5em}{}
\titlespacing*{\subsubsection}{0pt}{2.2ex plus .8ex minus .2ex}{1ex plus .2ex}
\titleformat{name=\subsubsection,numberless}[block]
  {\normalfont\normalsize\headfont}{}{0em}{}
\setcounter{secnumdepth}{3}
\setcounter{tocdepth}{3}
% 为含长外文书名与行内公式的段落提供必要的断行余量
\setlength{\emergencystretch}{3em}

% 章标题样式: 原书章标题与节同级排印("13. A) HERITAGE..."),
% 故章标题取块状左对齐黑体, 字号介于节与部之间
\titleformat{\chapter}[block]
  {\normalfont\LARGE\headfont}{\thechapter.}{0.5em}{}
\titlespacing*{\chapter}{0pt}{3.6ex plus 1.4ex minus .2ex}{2ex plus .2ex}
% 目录中章号照原书用阿拉伯数字("13. A) ..." 而非 "第十三章")
\renewcommand{\CTEXthechapter}{\thechapter\hspace{.3em}}

% --------------------------------------------------------------
%  目录观感: 条目与引导点之间、引导点与页码之间各留 0.5em
%  内核 \@dottedtocline 让 \leaders 从标题末字起排, 第一个点几乎
%  贴着末字, 读起来像多标了一个句点; 页码前也没有余量。两处各补
%  0.5em。
% --------------------------------------------------------------
\makeatletter
\patchcmd{\@dottedtocline}{\nobreak\leaders}{\nobreak\hspace{0.5em}\leaders}%
  {\typeout{MT: TOC 引导点前已留 0.5em}}%
  {\typeout{MT-WARN: TOC 引导点前间距补丁未生效}}
\patchcmd{\@dottedtocline}{\hfill\nobreak\hb@xt@\@pnumwidth}%
  {\hfill\nobreak\hspace{0.5em}\hb@xt@\@pnumwidth}%
  {\typeout{MT: TOC 页码前已留 0.5em}}%
  {\typeout{MT-WARN: TOC 页码前间距补丁未生效}}
\makeatother

% --------------------------------------------------------------
%  正文半角标点后的间距
%  本书体例用半角 , ; : ? ! 与 ASCII 括号, 由 xeCJK 按 HalfRight
%  (右半留白) 处理。放宽到 0.45em, 与全角标点右半留白大致齐平。
% --------------------------------------------------------------
\xeCJKsetup{CJKecglue={\hskip 0.45em plus 0.06em minus 0.09em}}

% --------------------------------------------------------------
%  页眉页脚: 偶数页眉右对齐斜体节题; 页码页脚外侧
% --------------------------------------------------------------
\fancyhf{}
\fancyhead[RE]{\small\itshape\rightmark}
\fancyfoot[LE,RO]{\small\thepage}
\renewcommand{\headrulewidth}{0pt}
\renewcommand{\sectionmark}[1]{\markright{\thesection.\ #1}}
\fancypagestyle{plain}{%
  \fancyhf{}%
  \fancyfoot[LE,RO]{\small \thepage}%
  \renewcommand{\headrulewidth}{0pt}%
}

% --------------------------------------------------------------
%  原书版式记号宏
% --------------------------------------------------------------
% 边码: 原打字稿页码(G. 手稿系统里的自引页码), 排在页边
\usepackage{marginnote}
\newcommand{\origpage}[1]{\marginnote{\normalfont\scriptsize p.~#1}}

% ⋄ 记号: 段首/行内正常大小; 上标形式为作者后补文字的位置标记
\DeclareRobustCommand{\pmark}{\ensuremath{\diamond}\hspace{0.15em}}
\DeclareRobustCommand{\dmark}{\textsuperscript{\ensuremath{\diamond}}}

% 脚注人工编号: \fnote{795}{(*)}{文本} —— 编号与(*)角标均照录原书
\newcommand{\fnote}[3]{%
  \begingroup
  \renewcommand{\thefootnote}{#1#2}%
  \footnote{#3}%
  \endgroup
}

% 第一部分的脚注: 页内序号作上标(#1), 注文行首印 "序号(连续注号)",
% 正文在序号之后另行括注同一连续注号 —— \fnoteP{1}{(3)}{注文}
% 原书有三条注无连续注号而改标日期(ch07 注3 = 30 Septembre,
% ch09 注11/12 = 8 août), 第二参数留空, 日期写进注文之首。
\newcommand{\fnoteP}[3]{%
  \begingroup
  \renewcommand{\thefootnote}{#1}\footnotemark
  \renewcommand{\thefootnote}{#1#2}\footnotetext{#3}%
  \endgroup
}

% 正文呼号与注文行首标号不同(原书照录)
\newcommand{\fnotevariant}[4]{%
  \begingroup
  \renewcommand{\thefootnote}{#1#2}\footnotemark
  \renewcommand{\thefootnote}{#1#3}\footnotetext{#4}%
  \endgroup
}

% 大注块标题: \notehead{171³}{斜体评论} —— 正文随后另起
\newcommand{\notehead}[2]{%
  \par\bigskip
  \noindent{\headfont 注~#1}\ifblank{#2}{}{\hspace{0.5em}{\itshape #2}}%
  \par\nobreak\smallskip\nobreak
}

% ∗ ∗ ∗ 分隔
\newcommand{\asterisms}{%
  \par\vspace{1.2ex}%
  \noindent\hspace*{2em}$\ast$\hspace{1.2em}$\ast$\hspace{1.2em}$\ast$%
  \par\vspace{1.2ex}%
}

% 无编号子目的目录条目
\newcommand{\tocsubsub}[1]{%
  \phantomsection
  \addcontentsline{toc}{subsubsection}{#1}%
}

% 部扉: 仿原书部题页
\newcommand{\partdivider}[3]{%
  \clearpage
  \thispagestyle{empty}%
  \phantomsection
  \addcontentsline{toc}{part}{#1\quad #3}%
  \vspace*{6.5cm}
  \begin{center}
    {\fontsize{16pt}{26pt}\selectfont\headfont #1\par}
    \vspace{0.9cm}
    {\fontsize{22pt}{34pt}\selectfont\headfont #2\par}
  \end{center}
  \clearpage
}

% 献辞页: 仿原书部题后的题献, 独占一页且不入目录
\newcommand{\dedication}[1]{%
  \thispagestyle{empty}%
  \vspace*{4cm}%
  \begin{center}#1\end{center}%
  \clearpage
}

% --------------------------------------------------------------
%  各卷"译者说明"的共用体例段(各卷的范围段自写, 后接本段)
% --------------------------------------------------------------
\newcommand{\CommonNoteBody}{%
为便于与原书及学界通行的打字稿页码对勘, 本册保留:

\begin{itemize}
  \item 原书全部章节编号(章~1--18, 节如~5.2、18.2, 子目如~18.1.1)与原书小节标题;
  \item 原书脚注的编号与角标——
        \begin{itemize}
          \item 前奏(第~1--4 章)与第一部的注在各页内重新编号, 正文中的上标即该
                页序号; 第一部紧随其后的括注(如“(3)”)是该部分的连续注号, 与第~12
                章“第一部分的注”里的“注~3”一一对应, 注文行首照原书印作“1(3)”式。
                其中三条注原书未标连续注号而改标日期(ch7 注~3 = 30 Septembre、
                ch9 注~11/12 = 8 août), 本册照录其日期;
          \item 第二部第~13--16~章各自重新编号, 分别为~1--41、1--86、1--132、1--29;
                第三部第~17~章自~1~起至~12, 第~18~章自~1~重新起算并延续至~794;
                末段自~\textbf{795(*)}~续至~1053(*)。各段与原书一一对应;
        \end{itemize}
  \item 页边码“p.~1”“p.~A3”“p.~1066”式小注——这是作者引用自己打字稿原稿时
        所用的页码, 与本排印本页码不同。前奏四章的边码各有前缀: 第~1 章作
        \texttt{A1…A6}, 第~2 章作 \texttt{P1…P65}, 第~3 章作 \texttt{L1…L56},
        第~4 章作罗马数字 \texttt{I…XXII}, 自第一部起作纯数字 \texttt{1…171};
  \item 作者用以标记段落的~$\diamond$~记号, 以及标示后来补写文字的上标
        ~\textsuperscript{$\diamond$}~记号;
  \item 原书大小注(“注~44'”“注~46$^1$”“注~171³”等)的编号与层级。
\end{itemize}

第~12 章“第一部分的注”把第~5 至第~11 章各页脚注汇为一编, 并给每条注配上节题
(如“12.19. 耶稣与十二使徒”), 因此部分注文与正文脚注重复——这是原书的体例, 非
重出之误。这一部分属随笔式的自省文字, 同一句法语在正文脚注与章末汇编注中
各有措辞, 本册有意保留各自的译笔, 不强求统一。

除上述标记外, 译文力求贴近作者驳杂、跳跃而辛辣的文风; 数学名词的译法见各卷
术语表(如~\emph{topos}~不译, \emph{faisceaux pervers}~作“反常层”)。文中的日期
括注(如“(5~月~30~日)”“(1984年3月补写)”)为作者历次补写时的纪日, 一律照录。文中
另有一些对注号、页码的引用(如“n°~1695”“打字稿第~227~页”)沿袭作者打字稿的编号
系统, 与本排印本不尽一致, 亦无法在本书中落实, 悉照原文录出。

全套五卷(另有全书合订一册)的分工是:

\begin{itemize}
  \item \textbf{卷一·前奏}——主题的呈现: 或四个乐章的前奏(原书 p.23--122);
  \item \textbf{卷二·第一部}——愚妄与更新(原书 p.123--254);
  \item \textbf{卷三·第二部}——葬礼(I): 或中国皇帝的新衣(原书 p.255--412);
  \item \textbf{卷四·第三部}——葬礼(II): 或阴阳之钥(原书 p.413--632);
  \item \textbf{卷五·第四部}——葬礼(III): 或四种运算(原书 p.633--929)。
\end{itemize}
}
 之后可用 \hypersetup 覆写)
% --------------------------------------------------------------
\hypersetup{
  pdftitle={收获与播种},
  pdfauthor={Alexandre Grothendieck / 数译 译},
  pdfsubject={Récoltes et Semailles — Réflexions et témoignage sur un passé de mathématicien}
}

% --------------------------------------------------------------
%  标题样式: 仿原书 —— 黑体左对齐, 编号后带句点
% --------------------------------------------------------------
\titleformat{\section}[block]
  {\normalfont\Large\headfont}{\thesection.}{0.5em}{}
\titlespacing*{\section}{0pt}{3.2ex plus 1.2ex minus .2ex}{1.6ex plus .2ex}
\titleformat{\subsection}[block]
  {\normalfont\large\headfont}{\thesubsection.}{0.5em}{}
\titlespacing*{\subsection}{0pt}{2.6ex plus 1ex minus .2ex}{1.2ex plus .2ex}
\titleformat{\subsubsection}[block]
  {\normalfont\normalsize\headfont}{\thesubsubsection.}{0.5em}{}
\titlespacing*{\subsubsection}{0pt}{2.2ex plus .8ex minus .2ex}{1ex plus .2ex}
\titleformat{name=\subsubsection,numberless}[block]
  {\normalfont\normalsize\headfont}{}{0em}{}
\setcounter{secnumdepth}{3}
\setcounter{tocdepth}{3}
% 为含长外文书名与行内公式的段落提供必要的断行余量
\setlength{\emergencystretch}{3em}

% 章标题样式: 原书章标题与节同级排印("13. A) HERITAGE..."),
% 故章标题取块状左对齐黑体, 字号介于节与部之间
\titleformat{\chapter}[block]
  {\normalfont\LARGE\headfont}{\thechapter.}{0.5em}{}
\titlespacing*{\chapter}{0pt}{3.6ex plus 1.4ex minus .2ex}{2ex plus .2ex}
% 目录中章号照原书用阿拉伯数字("13. A) ..." 而非 "第十三章")
\renewcommand{\CTEXthechapter}{\thechapter\hspace{.3em}}

% --------------------------------------------------------------
%  目录观感: 条目与引导点之间、引导点与页码之间各留 0.5em
%  内核 \@dottedtocline 让 \leaders 从标题末字起排, 第一个点几乎
%  贴着末字, 读起来像多标了一个句点; 页码前也没有余量。两处各补
%  0.5em。
% --------------------------------------------------------------
\makeatletter
\patchcmd{\@dottedtocline}{\nobreak\leaders}{\nobreak\hspace{0.5em}\leaders}%
  {\typeout{MT: TOC 引导点前已留 0.5em}}%
  {\typeout{MT-WARN: TOC 引导点前间距补丁未生效}}
\patchcmd{\@dottedtocline}{\hfill\nobreak\hb@xt@\@pnumwidth}%
  {\hfill\nobreak\hspace{0.5em}\hb@xt@\@pnumwidth}%
  {\typeout{MT: TOC 页码前已留 0.5em}}%
  {\typeout{MT-WARN: TOC 页码前间距补丁未生效}}
\makeatother

% --------------------------------------------------------------
%  正文半角标点后的间距
%  本书体例用半角 , ; : ? ! 与 ASCII 括号, 由 xeCJK 按 HalfRight
%  (右半留白) 处理。放宽到 0.45em, 与全角标点右半留白大致齐平。
% --------------------------------------------------------------
\xeCJKsetup{CJKecglue={\hskip 0.45em plus 0.06em minus 0.09em}}

% --------------------------------------------------------------
%  页眉页脚: 偶数页眉右对齐斜体节题; 页码页脚外侧
% --------------------------------------------------------------
\fancyhf{}
\fancyhead[RE]{\small\itshape\rightmark}
\fancyfoot[LE,RO]{\small\thepage}
\renewcommand{\headrulewidth}{0pt}
\renewcommand{\sectionmark}[1]{\markright{\thesection.\ #1}}
\fancypagestyle{plain}{%
  \fancyhf{}%
  \fancyfoot[LE,RO]{\small \thepage}%
  \renewcommand{\headrulewidth}{0pt}%
}

% --------------------------------------------------------------
%  原书版式记号宏
% --------------------------------------------------------------
% 边码: 原打字稿页码(G. 手稿系统里的自引页码), 排在页边
\usepackage{marginnote}
\newcommand{\origpage}[1]{\marginnote{\normalfont\scriptsize p.~#1}}

% ⋄ 记号: 段首/行内正常大小; 上标形式为作者后补文字的位置标记
\DeclareRobustCommand{\pmark}{\ensuremath{\diamond}\hspace{0.15em}}
\DeclareRobustCommand{\dmark}{\textsuperscript{\ensuremath{\diamond}}}

% 脚注人工编号: \fnote{795}{(*)}{文本} —— 编号与(*)角标均照录原书
\newcommand{\fnote}[3]{%
  \begingroup
  \renewcommand{\thefootnote}{#1#2}%
  \footnote{#3}%
  \endgroup
}

% 第一部分的脚注: 页内序号作上标(#1), 注文行首印 "序号(连续注号)",
% 正文在序号之后另行括注同一连续注号 —— \fnoteP{1}{(3)}{注文}
% 原书有三条注无连续注号而改标日期(ch07 注3 = 30 Septembre,
% ch09 注11/12 = 8 août), 第二参数留空, 日期写进注文之首。
\newcommand{\fnoteP}[3]{%
  \begingroup
  \renewcommand{\thefootnote}{#1}\footnotemark
  \renewcommand{\thefootnote}{#1#2}\footnotetext{#3}%
  \endgroup
}

% 正文呼号与注文行首标号不同(原书照录)
\newcommand{\fnotevariant}[4]{%
  \begingroup
  \renewcommand{\thefootnote}{#1#2}\footnotemark
  \renewcommand{\thefootnote}{#1#3}\footnotetext{#4}%
  \endgroup
}

% 大注块标题: \notehead{171³}{斜体评论} —— 正文随后另起
\newcommand{\notehead}[2]{%
  \par\bigskip
  \noindent{\headfont 注~#1}\ifblank{#2}{}{\hspace{0.5em}{\itshape #2}}%
  \par\nobreak\smallskip\nobreak
}

% ∗ ∗ ∗ 分隔
\newcommand{\asterisms}{%
  \par\vspace{1.2ex}%
  \noindent\hspace*{2em}$\ast$\hspace{1.2em}$\ast$\hspace{1.2em}$\ast$%
  \par\vspace{1.2ex}%
}

% 无编号子目的目录条目
\newcommand{\tocsubsub}[1]{%
  \phantomsection
  \addcontentsline{toc}{subsubsection}{#1}%
}

% 部扉: 仿原书部题页
\newcommand{\partdivider}[3]{%
  \clearpage
  \thispagestyle{empty}%
  \phantomsection
  \addcontentsline{toc}{part}{#1\quad #3}%
  \vspace*{6.5cm}
  \begin{center}
    {\fontsize{16pt}{26pt}\selectfont\headfont #1\par}
    \vspace{0.9cm}
    {\fontsize{22pt}{34pt}\selectfont\headfont #2\par}
  \end{center}
  \clearpage
}

% 献辞页: 仿原书部题后的题献, 独占一页且不入目录
\newcommand{\dedication}[1]{%
  \thispagestyle{empty}%
  \vspace*{4cm}%
  \begin{center}#1\end{center}%
  \clearpage
}

% --------------------------------------------------------------
%  各卷"译者说明"的共用体例段(各卷的范围段自写, 后接本段)
% --------------------------------------------------------------
\newcommand{\CommonNoteBody}{%
为便于与原书及学界通行的打字稿页码对勘, 本册保留:

\begin{itemize}
  \item 原书全部章节编号(章~1--18, 节如~5.2、18.2, 子目如~18.1.1)与原书小节标题;
  \item 原书脚注的编号与角标——
        \begin{itemize}
          \item 前奏(第~1--4 章)与第一部的注在各页内重新编号, 正文中的上标即该
                页序号; 第一部紧随其后的括注(如“(3)”)是该部分的连续注号, 与第~12
                章“第一部分的注”里的“注~3”一一对应, 注文行首照原书印作“1(3)”式。
                其中三条注原书未标连续注号而改标日期(ch7 注~3 = 30 Septembre、
                ch9 注~11/12 = 8 août), 本册照录其日期;
          \item 第二部第~13--16~章各自重新编号, 分别为~1--41、1--86、1--132、1--29;
                第三部第~17~章自~1~起至~12, 第~18~章自~1~重新起算并延续至~794;
                末段自~\textbf{795(*)}~续至~1053(*)。各段与原书一一对应;
        \end{itemize}
  \item 页边码“p.~1”“p.~A3”“p.~1066”式小注——这是作者引用自己打字稿原稿时
        所用的页码, 与本排印本页码不同。前奏四章的边码各有前缀: 第~1 章作
        \texttt{A1…A6}, 第~2 章作 \texttt{P1…P65}, 第~3 章作 \texttt{L1…L56},
        第~4 章作罗马数字 \texttt{I…XXII}, 自第一部起作纯数字 \texttt{1…171};
  \item 作者用以标记段落的~$\diamond$~记号, 以及标示后来补写文字的上标
        ~\textsuperscript{$\diamond$}~记号;
  \item 原书大小注(“注~44'”“注~46$^1$”“注~171³”等)的编号与层级。
\end{itemize}

第~12 章“第一部分的注”把第~5 至第~11 章各页脚注汇为一编, 并给每条注配上节题
(如“12.19. 耶稣与十二使徒”), 因此部分注文与正文脚注重复——这是原书的体例, 非
重出之误。这一部分属随笔式的自省文字, 同一句法语在正文脚注与章末汇编注中
各有措辞, 本册有意保留各自的译笔, 不强求统一。

除上述标记外, 译文力求贴近作者驳杂、跳跃而辛辣的文风; 数学名词的译法见各卷
术语表(如~\emph{topos}~不译, \emph{faisceaux pervers}~作“反常层”)。文中的日期
括注(如“(5~月~30~日)”“(1984年3月补写)”)为作者历次补写时的纪日, 一律照录。文中
另有一些对注号、页码的引用(如“n°~1695”“打字稿第~227~页”)沿袭作者打字稿的编号
系统, 与本排印本不尽一致, 亦无法在本书中落实, 悉照原文录出。

全套五卷(另有全书合订一册)的分工是:

\begin{itemize}
  \item \textbf{卷一·前奏}——主题的呈现: 或四个乐章的前奏(原书 p.23--122);
  \item \textbf{卷二·第一部}——愚妄与更新(原书 p.123--254);
  \item \textbf{卷三·第二部}——葬礼(I): 或中国皇帝的新衣(原书 p.255--412);
  \item \textbf{卷四·第三部}——葬礼(II): 或阴阳之钥(原书 p.413--632);
  \item \textbf{卷五·第四部}——葬礼(III): 或四种运算(原书 p.633--929)。
\end{itemize}
}
                       % 共用导言区
%      % ================================================================
%  《收获与播种》统一封面 —— TikZ 原生矢量绘制(不依赖任何外部图片)
%  母题: 暮色中的麦田 —— 犁沟向地平线的灭点收束, 前景麦穗剪影,
%        一位播种者把一把种子撒向天空, 种子沿弧线落入田里。
%  用法(在 \begin{document} 之前, 依次):
%      % ================================================================
%  《收获与播种》中文译本 —— 各卷共用导言区
%  由 vol-prelude / vol-part1 / vol-part2 / vol-part3 / vol-part4
%  以及全书合订 main-all 共同 \input。基类 mathtranslation.cls v3.2 不改。
%  --------------------------------------------------------------
%  样式均在此改写, 以贴近原书:
%    · 章标题仿原书 "5. Travail et découverte"(块状左对齐黑体)
%    · 节/小节: 黑体左对齐, 编号后带句点
%    · 脚注: 人工编号照录原书——前奏与第一部分页内重新编号,
%      第二至四部各自连续编号并带(*)角标; 第一部分另有 \fnoteP
%      处理 "(n)" 式连续注号
%    · 边码: 原打字稿页码 \origpage{N}(前奏作 A1、P1、L1、罗马数字等)
%    · ⋄ 记号: 段首补写标记 \pmark 与上标形式 \dmark
% ================================================================

% 题名含空格, 不能走类选项(会被 kvoptions 吃掉), 在此赋值
\renewcommand{\BookTitleEN}{Récoltes et Semailles}
\renewcommand{\OriginalAuthor}{Alexandre Grothendieck}
\renewcommand{\TranslationDate}{\today}

% --------------------------------------------------------------
%  罕用字后备(如"祐"): Fandol 不载时退到 SimSun
% --------------------------------------------------------------
\setCJKfallbackfamilyfont{\CJKrmdefault}{{SimSun}}

% --------------------------------------------------------------
%  hyperref 元信息(各卷在 % ================================================================
%  《收获与播种》中文译本 —— 各卷共用导言区
%  由 vol-prelude / vol-part1 / vol-part2 / vol-part3 / vol-part4
%  以及全书合订 main-all 共同 \input。基类 mathtranslation.cls v3.2 不改。
%  --------------------------------------------------------------
%  样式均在此改写, 以贴近原书:
%    · 章标题仿原书 "5. Travail et découverte"(块状左对齐黑体)
%    · 节/小节: 黑体左对齐, 编号后带句点
%    · 脚注: 人工编号照录原书——前奏与第一部分页内重新编号,
%      第二至四部各自连续编号并带(*)角标; 第一部分另有 \fnoteP
%      处理 "(n)" 式连续注号
%    · 边码: 原打字稿页码 \origpage{N}(前奏作 A1、P1、L1、罗马数字等)
%    · ⋄ 记号: 段首补写标记 \pmark 与上标形式 \dmark
% ================================================================

% 题名含空格, 不能走类选项(会被 kvoptions 吃掉), 在此赋值
\renewcommand{\BookTitleEN}{Récoltes et Semailles}
\renewcommand{\OriginalAuthor}{Alexandre Grothendieck}
\renewcommand{\TranslationDate}{\today}

% --------------------------------------------------------------
%  罕用字后备(如"祐"): Fandol 不载时退到 SimSun
% --------------------------------------------------------------
\setCJKfallbackfamilyfont{\CJKrmdefault}{{SimSun}}

% --------------------------------------------------------------
%  hyperref 元信息(各卷在 \input{common} 之后可用 \hypersetup 覆写)
% --------------------------------------------------------------
\hypersetup{
  pdftitle={收获与播种},
  pdfauthor={Alexandre Grothendieck / 数译 译},
  pdfsubject={Récoltes et Semailles — Réflexions et témoignage sur un passé de mathématicien}
}

% --------------------------------------------------------------
%  标题样式: 仿原书 —— 黑体左对齐, 编号后带句点
% --------------------------------------------------------------
\titleformat{\section}[block]
  {\normalfont\Large\headfont}{\thesection.}{0.5em}{}
\titlespacing*{\section}{0pt}{3.2ex plus 1.2ex minus .2ex}{1.6ex plus .2ex}
\titleformat{\subsection}[block]
  {\normalfont\large\headfont}{\thesubsection.}{0.5em}{}
\titlespacing*{\subsection}{0pt}{2.6ex plus 1ex minus .2ex}{1.2ex plus .2ex}
\titleformat{\subsubsection}[block]
  {\normalfont\normalsize\headfont}{\thesubsubsection.}{0.5em}{}
\titlespacing*{\subsubsection}{0pt}{2.2ex plus .8ex minus .2ex}{1ex plus .2ex}
\titleformat{name=\subsubsection,numberless}[block]
  {\normalfont\normalsize\headfont}{}{0em}{}
\setcounter{secnumdepth}{3}
\setcounter{tocdepth}{3}
% 为含长外文书名与行内公式的段落提供必要的断行余量
\setlength{\emergencystretch}{3em}

% 章标题样式: 原书章标题与节同级排印("13. A) HERITAGE..."),
% 故章标题取块状左对齐黑体, 字号介于节与部之间
\titleformat{\chapter}[block]
  {\normalfont\LARGE\headfont}{\thechapter.}{0.5em}{}
\titlespacing*{\chapter}{0pt}{3.6ex plus 1.4ex minus .2ex}{2ex plus .2ex}
% 目录中章号照原书用阿拉伯数字("13. A) ..." 而非 "第十三章")
\renewcommand{\CTEXthechapter}{\thechapter\hspace{.3em}}

% --------------------------------------------------------------
%  目录观感: 条目与引导点之间、引导点与页码之间各留 0.5em
%  内核 \@dottedtocline 让 \leaders 从标题末字起排, 第一个点几乎
%  贴着末字, 读起来像多标了一个句点; 页码前也没有余量。两处各补
%  0.5em。
% --------------------------------------------------------------
\makeatletter
\patchcmd{\@dottedtocline}{\nobreak\leaders}{\nobreak\hspace{0.5em}\leaders}%
  {\typeout{MT: TOC 引导点前已留 0.5em}}%
  {\typeout{MT-WARN: TOC 引导点前间距补丁未生效}}
\patchcmd{\@dottedtocline}{\hfill\nobreak\hb@xt@\@pnumwidth}%
  {\hfill\nobreak\hspace{0.5em}\hb@xt@\@pnumwidth}%
  {\typeout{MT: TOC 页码前已留 0.5em}}%
  {\typeout{MT-WARN: TOC 页码前间距补丁未生效}}
\makeatother

% --------------------------------------------------------------
%  正文半角标点后的间距
%  本书体例用半角 , ; : ? ! 与 ASCII 括号, 由 xeCJK 按 HalfRight
%  (右半留白) 处理。放宽到 0.45em, 与全角标点右半留白大致齐平。
% --------------------------------------------------------------
\xeCJKsetup{CJKecglue={\hskip 0.45em plus 0.06em minus 0.09em}}

% --------------------------------------------------------------
%  页眉页脚: 偶数页眉右对齐斜体节题; 页码页脚外侧
% --------------------------------------------------------------
\fancyhf{}
\fancyhead[RE]{\small\itshape\rightmark}
\fancyfoot[LE,RO]{\small\thepage}
\renewcommand{\headrulewidth}{0pt}
\renewcommand{\sectionmark}[1]{\markright{\thesection.\ #1}}
\fancypagestyle{plain}{%
  \fancyhf{}%
  \fancyfoot[LE,RO]{\small \thepage}%
  \renewcommand{\headrulewidth}{0pt}%
}

% --------------------------------------------------------------
%  原书版式记号宏
% --------------------------------------------------------------
% 边码: 原打字稿页码(G. 手稿系统里的自引页码), 排在页边
\usepackage{marginnote}
\newcommand{\origpage}[1]{\marginnote{\normalfont\scriptsize p.~#1}}

% ⋄ 记号: 段首/行内正常大小; 上标形式为作者后补文字的位置标记
\DeclareRobustCommand{\pmark}{\ensuremath{\diamond}\hspace{0.15em}}
\DeclareRobustCommand{\dmark}{\textsuperscript{\ensuremath{\diamond}}}

% 脚注人工编号: \fnote{795}{(*)}{文本} —— 编号与(*)角标均照录原书
\newcommand{\fnote}[3]{%
  \begingroup
  \renewcommand{\thefootnote}{#1#2}%
  \footnote{#3}%
  \endgroup
}

% 第一部分的脚注: 页内序号作上标(#1), 注文行首印 "序号(连续注号)",
% 正文在序号之后另行括注同一连续注号 —— \fnoteP{1}{(3)}{注文}
% 原书有三条注无连续注号而改标日期(ch07 注3 = 30 Septembre,
% ch09 注11/12 = 8 août), 第二参数留空, 日期写进注文之首。
\newcommand{\fnoteP}[3]{%
  \begingroup
  \renewcommand{\thefootnote}{#1}\footnotemark
  \renewcommand{\thefootnote}{#1#2}\footnotetext{#3}%
  \endgroup
}

% 正文呼号与注文行首标号不同(原书照录)
\newcommand{\fnotevariant}[4]{%
  \begingroup
  \renewcommand{\thefootnote}{#1#2}\footnotemark
  \renewcommand{\thefootnote}{#1#3}\footnotetext{#4}%
  \endgroup
}

% 大注块标题: \notehead{171³}{斜体评论} —— 正文随后另起
\newcommand{\notehead}[2]{%
  \par\bigskip
  \noindent{\headfont 注~#1}\ifblank{#2}{}{\hspace{0.5em}{\itshape #2}}%
  \par\nobreak\smallskip\nobreak
}

% ∗ ∗ ∗ 分隔
\newcommand{\asterisms}{%
  \par\vspace{1.2ex}%
  \noindent\hspace*{2em}$\ast$\hspace{1.2em}$\ast$\hspace{1.2em}$\ast$%
  \par\vspace{1.2ex}%
}

% 无编号子目的目录条目
\newcommand{\tocsubsub}[1]{%
  \phantomsection
  \addcontentsline{toc}{subsubsection}{#1}%
}

% 部扉: 仿原书部题页
\newcommand{\partdivider}[3]{%
  \clearpage
  \thispagestyle{empty}%
  \phantomsection
  \addcontentsline{toc}{part}{#1\quad #3}%
  \vspace*{6.5cm}
  \begin{center}
    {\fontsize{16pt}{26pt}\selectfont\headfont #1\par}
    \vspace{0.9cm}
    {\fontsize{22pt}{34pt}\selectfont\headfont #2\par}
  \end{center}
  \clearpage
}

% 献辞页: 仿原书部题后的题献, 独占一页且不入目录
\newcommand{\dedication}[1]{%
  \thispagestyle{empty}%
  \vspace*{4cm}%
  \begin{center}#1\end{center}%
  \clearpage
}

% --------------------------------------------------------------
%  各卷"译者说明"的共用体例段(各卷的范围段自写, 后接本段)
% --------------------------------------------------------------
\newcommand{\CommonNoteBody}{%
为便于与原书及学界通行的打字稿页码对勘, 本册保留:

\begin{itemize}
  \item 原书全部章节编号(章~1--18, 节如~5.2、18.2, 子目如~18.1.1)与原书小节标题;
  \item 原书脚注的编号与角标——
        \begin{itemize}
          \item 前奏(第~1--4 章)与第一部的注在各页内重新编号, 正文中的上标即该
                页序号; 第一部紧随其后的括注(如“(3)”)是该部分的连续注号, 与第~12
                章“第一部分的注”里的“注~3”一一对应, 注文行首照原书印作“1(3)”式。
                其中三条注原书未标连续注号而改标日期(ch7 注~3 = 30 Septembre、
                ch9 注~11/12 = 8 août), 本册照录其日期;
          \item 第二部第~13--16~章各自重新编号, 分别为~1--41、1--86、1--132、1--29;
                第三部第~17~章自~1~起至~12, 第~18~章自~1~重新起算并延续至~794;
                末段自~\textbf{795(*)}~续至~1053(*)。各段与原书一一对应;
        \end{itemize}
  \item 页边码“p.~1”“p.~A3”“p.~1066”式小注——这是作者引用自己打字稿原稿时
        所用的页码, 与本排印本页码不同。前奏四章的边码各有前缀: 第~1 章作
        \texttt{A1…A6}, 第~2 章作 \texttt{P1…P65}, 第~3 章作 \texttt{L1…L56},
        第~4 章作罗马数字 \texttt{I…XXII}, 自第一部起作纯数字 \texttt{1…171};
  \item 作者用以标记段落的~$\diamond$~记号, 以及标示后来补写文字的上标
        ~\textsuperscript{$\diamond$}~记号;
  \item 原书大小注(“注~44'”“注~46$^1$”“注~171³”等)的编号与层级。
\end{itemize}

第~12 章“第一部分的注”把第~5 至第~11 章各页脚注汇为一编, 并给每条注配上节题
(如“12.19. 耶稣与十二使徒”), 因此部分注文与正文脚注重复——这是原书的体例, 非
重出之误。这一部分属随笔式的自省文字, 同一句法语在正文脚注与章末汇编注中
各有措辞, 本册有意保留各自的译笔, 不强求统一。

除上述标记外, 译文力求贴近作者驳杂、跳跃而辛辣的文风; 数学名词的译法见各卷
术语表(如~\emph{topos}~不译, \emph{faisceaux pervers}~作“反常层”)。文中的日期
括注(如“(5~月~30~日)”“(1984年3月补写)”)为作者历次补写时的纪日, 一律照录。文中
另有一些对注号、页码的引用(如“n°~1695”“打字稿第~227~页”)沿袭作者打字稿的编号
系统, 与本排印本不尽一致, 亦无法在本书中落实, 悉照原文录出。

全套五卷(另有全书合订一册)的分工是:

\begin{itemize}
  \item \textbf{卷一·前奏}——主题的呈现: 或四个乐章的前奏(原书 p.23--122);
  \item \textbf{卷二·第一部}——愚妄与更新(原书 p.123--254);
  \item \textbf{卷三·第二部}——葬礼(I): 或中国皇帝的新衣(原书 p.255--412);
  \item \textbf{卷四·第三部}——葬礼(II): 或阴阳之钥(原书 p.413--632);
  \item \textbf{卷五·第四部}——葬礼(III): 或四种运算(原书 p.633--929)。
\end{itemize}
}
 之后可用 \hypersetup 覆写)
% --------------------------------------------------------------
\hypersetup{
  pdftitle={收获与播种},
  pdfauthor={Alexandre Grothendieck / 数译 译},
  pdfsubject={Récoltes et Semailles — Réflexions et témoignage sur un passé de mathématicien}
}

% --------------------------------------------------------------
%  标题样式: 仿原书 —— 黑体左对齐, 编号后带句点
% --------------------------------------------------------------
\titleformat{\section}[block]
  {\normalfont\Large\headfont}{\thesection.}{0.5em}{}
\titlespacing*{\section}{0pt}{3.2ex plus 1.2ex minus .2ex}{1.6ex plus .2ex}
\titleformat{\subsection}[block]
  {\normalfont\large\headfont}{\thesubsection.}{0.5em}{}
\titlespacing*{\subsection}{0pt}{2.6ex plus 1ex minus .2ex}{1.2ex plus .2ex}
\titleformat{\subsubsection}[block]
  {\normalfont\normalsize\headfont}{\thesubsubsection.}{0.5em}{}
\titlespacing*{\subsubsection}{0pt}{2.2ex plus .8ex minus .2ex}{1ex plus .2ex}
\titleformat{name=\subsubsection,numberless}[block]
  {\normalfont\normalsize\headfont}{}{0em}{}
\setcounter{secnumdepth}{3}
\setcounter{tocdepth}{3}
% 为含长外文书名与行内公式的段落提供必要的断行余量
\setlength{\emergencystretch}{3em}

% 章标题样式: 原书章标题与节同级排印("13. A) HERITAGE..."),
% 故章标题取块状左对齐黑体, 字号介于节与部之间
\titleformat{\chapter}[block]
  {\normalfont\LARGE\headfont}{\thechapter.}{0.5em}{}
\titlespacing*{\chapter}{0pt}{3.6ex plus 1.4ex minus .2ex}{2ex plus .2ex}
% 目录中章号照原书用阿拉伯数字("13. A) ..." 而非 "第十三章")
\renewcommand{\CTEXthechapter}{\thechapter\hspace{.3em}}

% --------------------------------------------------------------
%  目录观感: 条目与引导点之间、引导点与页码之间各留 0.5em
%  内核 \@dottedtocline 让 \leaders 从标题末字起排, 第一个点几乎
%  贴着末字, 读起来像多标了一个句点; 页码前也没有余量。两处各补
%  0.5em。
% --------------------------------------------------------------
\makeatletter
\patchcmd{\@dottedtocline}{\nobreak\leaders}{\nobreak\hspace{0.5em}\leaders}%
  {\typeout{MT: TOC 引导点前已留 0.5em}}%
  {\typeout{MT-WARN: TOC 引导点前间距补丁未生效}}
\patchcmd{\@dottedtocline}{\hfill\nobreak\hb@xt@\@pnumwidth}%
  {\hfill\nobreak\hspace{0.5em}\hb@xt@\@pnumwidth}%
  {\typeout{MT: TOC 页码前已留 0.5em}}%
  {\typeout{MT-WARN: TOC 页码前间距补丁未生效}}
\makeatother

% --------------------------------------------------------------
%  正文半角标点后的间距
%  本书体例用半角 , ; : ? ! 与 ASCII 括号, 由 xeCJK 按 HalfRight
%  (右半留白) 处理。放宽到 0.45em, 与全角标点右半留白大致齐平。
% --------------------------------------------------------------
\xeCJKsetup{CJKecglue={\hskip 0.45em plus 0.06em minus 0.09em}}

% --------------------------------------------------------------
%  页眉页脚: 偶数页眉右对齐斜体节题; 页码页脚外侧
% --------------------------------------------------------------
\fancyhf{}
\fancyhead[RE]{\small\itshape\rightmark}
\fancyfoot[LE,RO]{\small\thepage}
\renewcommand{\headrulewidth}{0pt}
\renewcommand{\sectionmark}[1]{\markright{\thesection.\ #1}}
\fancypagestyle{plain}{%
  \fancyhf{}%
  \fancyfoot[LE,RO]{\small \thepage}%
  \renewcommand{\headrulewidth}{0pt}%
}

% --------------------------------------------------------------
%  原书版式记号宏
% --------------------------------------------------------------
% 边码: 原打字稿页码(G. 手稿系统里的自引页码), 排在页边
\usepackage{marginnote}
\newcommand{\origpage}[1]{\marginnote{\normalfont\scriptsize p.~#1}}

% ⋄ 记号: 段首/行内正常大小; 上标形式为作者后补文字的位置标记
\DeclareRobustCommand{\pmark}{\ensuremath{\diamond}\hspace{0.15em}}
\DeclareRobustCommand{\dmark}{\textsuperscript{\ensuremath{\diamond}}}

% 脚注人工编号: \fnote{795}{(*)}{文本} —— 编号与(*)角标均照录原书
\newcommand{\fnote}[3]{%
  \begingroup
  \renewcommand{\thefootnote}{#1#2}%
  \footnote{#3}%
  \endgroup
}

% 第一部分的脚注: 页内序号作上标(#1), 注文行首印 "序号(连续注号)",
% 正文在序号之后另行括注同一连续注号 —— \fnoteP{1}{(3)}{注文}
% 原书有三条注无连续注号而改标日期(ch07 注3 = 30 Septembre,
% ch09 注11/12 = 8 août), 第二参数留空, 日期写进注文之首。
\newcommand{\fnoteP}[3]{%
  \begingroup
  \renewcommand{\thefootnote}{#1}\footnotemark
  \renewcommand{\thefootnote}{#1#2}\footnotetext{#3}%
  \endgroup
}

% 正文呼号与注文行首标号不同(原书照录)
\newcommand{\fnotevariant}[4]{%
  \begingroup
  \renewcommand{\thefootnote}{#1#2}\footnotemark
  \renewcommand{\thefootnote}{#1#3}\footnotetext{#4}%
  \endgroup
}

% 大注块标题: \notehead{171³}{斜体评论} —— 正文随后另起
\newcommand{\notehead}[2]{%
  \par\bigskip
  \noindent{\headfont 注~#1}\ifblank{#2}{}{\hspace{0.5em}{\itshape #2}}%
  \par\nobreak\smallskip\nobreak
}

% ∗ ∗ ∗ 分隔
\newcommand{\asterisms}{%
  \par\vspace{1.2ex}%
  \noindent\hspace*{2em}$\ast$\hspace{1.2em}$\ast$\hspace{1.2em}$\ast$%
  \par\vspace{1.2ex}%
}

% 无编号子目的目录条目
\newcommand{\tocsubsub}[1]{%
  \phantomsection
  \addcontentsline{toc}{subsubsection}{#1}%
}

% 部扉: 仿原书部题页
\newcommand{\partdivider}[3]{%
  \clearpage
  \thispagestyle{empty}%
  \phantomsection
  \addcontentsline{toc}{part}{#1\quad #3}%
  \vspace*{6.5cm}
  \begin{center}
    {\fontsize{16pt}{26pt}\selectfont\headfont #1\par}
    \vspace{0.9cm}
    {\fontsize{22pt}{34pt}\selectfont\headfont #2\par}
  \end{center}
  \clearpage
}

% 献辞页: 仿原书部题后的题献, 独占一页且不入目录
\newcommand{\dedication}[1]{%
  \thispagestyle{empty}%
  \vspace*{4cm}%
  \begin{center}#1\end{center}%
  \clearpage
}

% --------------------------------------------------------------
%  各卷"译者说明"的共用体例段(各卷的范围段自写, 后接本段)
% --------------------------------------------------------------
\newcommand{\CommonNoteBody}{%
为便于与原书及学界通行的打字稿页码对勘, 本册保留:

\begin{itemize}
  \item 原书全部章节编号(章~1--18, 节如~5.2、18.2, 子目如~18.1.1)与原书小节标题;
  \item 原书脚注的编号与角标——
        \begin{itemize}
          \item 前奏(第~1--4 章)与第一部的注在各页内重新编号, 正文中的上标即该
                页序号; 第一部紧随其后的括注(如“(3)”)是该部分的连续注号, 与第~12
                章“第一部分的注”里的“注~3”一一对应, 注文行首照原书印作“1(3)”式。
                其中三条注原书未标连续注号而改标日期(ch7 注~3 = 30 Septembre、
                ch9 注~11/12 = 8 août), 本册照录其日期;
          \item 第二部第~13--16~章各自重新编号, 分别为~1--41、1--86、1--132、1--29;
                第三部第~17~章自~1~起至~12, 第~18~章自~1~重新起算并延续至~794;
                末段自~\textbf{795(*)}~续至~1053(*)。各段与原书一一对应;
        \end{itemize}
  \item 页边码“p.~1”“p.~A3”“p.~1066”式小注——这是作者引用自己打字稿原稿时
        所用的页码, 与本排印本页码不同。前奏四章的边码各有前缀: 第~1 章作
        \texttt{A1…A6}, 第~2 章作 \texttt{P1…P65}, 第~3 章作 \texttt{L1…L56},
        第~4 章作罗马数字 \texttt{I…XXII}, 自第一部起作纯数字 \texttt{1…171};
  \item 作者用以标记段落的~$\diamond$~记号, 以及标示后来补写文字的上标
        ~\textsuperscript{$\diamond$}~记号;
  \item 原书大小注(“注~44'”“注~46$^1$”“注~171³”等)的编号与层级。
\end{itemize}

第~12 章“第一部分的注”把第~5 至第~11 章各页脚注汇为一编, 并给每条注配上节题
(如“12.19. 耶稣与十二使徒”), 因此部分注文与正文脚注重复——这是原书的体例, 非
重出之误。这一部分属随笔式的自省文字, 同一句法语在正文脚注与章末汇编注中
各有措辞, 本册有意保留各自的译笔, 不强求统一。

除上述标记外, 译文力求贴近作者驳杂、跳跃而辛辣的文风; 数学名词的译法见各卷
术语表(如~\emph{topos}~不译, \emph{faisceaux pervers}~作“反常层”)。文中的日期
括注(如“(5~月~30~日)”“(1984年3月补写)”)为作者历次补写时的纪日, 一律照录。文中
另有一些对注号、页码的引用(如“n°~1695”“打字稿第~227~页”)沿袭作者打字稿的编号
系统, 与本排印本不尽一致, 亦无法在本书中落实, 悉照原文录出。

全套五卷(另有全书合订一册)的分工是:

\begin{itemize}
  \item \textbf{卷一·前奏}——主题的呈现: 或四个乐章的前奏(原书 p.23--122);
  \item \textbf{卷二·第一部}——愚妄与更新(原书 p.123--254);
  \item \textbf{卷三·第二部}——葬礼(I): 或中国皇帝的新衣(原书 p.255--412);
  \item \textbf{卷四·第三部}——葬礼(II): 或阴阳之钥(原书 p.413--632);
  \item \textbf{卷五·第四部}——葬礼(III): 或四种运算(原书 p.633--929)。
\end{itemize}
}
                       % 共用导言区
%      % ================================================================
%  《收获与播种》统一封面 —— TikZ 原生矢量绘制(不依赖任何外部图片)
%  母题: 暮色中的麦田 —— 犁沟向地平线的灭点收束, 前景麦穗剪影,
%        一位播种者把一把种子撒向天空, 种子沿弧线落入田里。
%  用法(在 \begin{document} 之前, 依次):
%      % ================================================================
%  《收获与播种》中文译本 —— 各卷共用导言区
%  由 vol-prelude / vol-part1 / vol-part2 / vol-part3 / vol-part4
%  以及全书合订 main-all 共同 \input。基类 mathtranslation.cls v3.2 不改。
%  --------------------------------------------------------------
%  样式均在此改写, 以贴近原书:
%    · 章标题仿原书 "5. Travail et découverte"(块状左对齐黑体)
%    · 节/小节: 黑体左对齐, 编号后带句点
%    · 脚注: 人工编号照录原书——前奏与第一部分页内重新编号,
%      第二至四部各自连续编号并带(*)角标; 第一部分另有 \fnoteP
%      处理 "(n)" 式连续注号
%    · 边码: 原打字稿页码 \origpage{N}(前奏作 A1、P1、L1、罗马数字等)
%    · ⋄ 记号: 段首补写标记 \pmark 与上标形式 \dmark
% ================================================================

% 题名含空格, 不能走类选项(会被 kvoptions 吃掉), 在此赋值
\renewcommand{\BookTitleEN}{Récoltes et Semailles}
\renewcommand{\OriginalAuthor}{Alexandre Grothendieck}
\renewcommand{\TranslationDate}{\today}

% --------------------------------------------------------------
%  罕用字后备(如"祐"): Fandol 不载时退到 SimSun
% --------------------------------------------------------------
\setCJKfallbackfamilyfont{\CJKrmdefault}{{SimSun}}

% --------------------------------------------------------------
%  hyperref 元信息(各卷在 \input{common} 之后可用 \hypersetup 覆写)
% --------------------------------------------------------------
\hypersetup{
  pdftitle={收获与播种},
  pdfauthor={Alexandre Grothendieck / 数译 译},
  pdfsubject={Récoltes et Semailles — Réflexions et témoignage sur un passé de mathématicien}
}

% --------------------------------------------------------------
%  标题样式: 仿原书 —— 黑体左对齐, 编号后带句点
% --------------------------------------------------------------
\titleformat{\section}[block]
  {\normalfont\Large\headfont}{\thesection.}{0.5em}{}
\titlespacing*{\section}{0pt}{3.2ex plus 1.2ex minus .2ex}{1.6ex plus .2ex}
\titleformat{\subsection}[block]
  {\normalfont\large\headfont}{\thesubsection.}{0.5em}{}
\titlespacing*{\subsection}{0pt}{2.6ex plus 1ex minus .2ex}{1.2ex plus .2ex}
\titleformat{\subsubsection}[block]
  {\normalfont\normalsize\headfont}{\thesubsubsection.}{0.5em}{}
\titlespacing*{\subsubsection}{0pt}{2.2ex plus .8ex minus .2ex}{1ex plus .2ex}
\titleformat{name=\subsubsection,numberless}[block]
  {\normalfont\normalsize\headfont}{}{0em}{}
\setcounter{secnumdepth}{3}
\setcounter{tocdepth}{3}
% 为含长外文书名与行内公式的段落提供必要的断行余量
\setlength{\emergencystretch}{3em}

% 章标题样式: 原书章标题与节同级排印("13. A) HERITAGE..."),
% 故章标题取块状左对齐黑体, 字号介于节与部之间
\titleformat{\chapter}[block]
  {\normalfont\LARGE\headfont}{\thechapter.}{0.5em}{}
\titlespacing*{\chapter}{0pt}{3.6ex plus 1.4ex minus .2ex}{2ex plus .2ex}
% 目录中章号照原书用阿拉伯数字("13. A) ..." 而非 "第十三章")
\renewcommand{\CTEXthechapter}{\thechapter\hspace{.3em}}

% --------------------------------------------------------------
%  目录观感: 条目与引导点之间、引导点与页码之间各留 0.5em
%  内核 \@dottedtocline 让 \leaders 从标题末字起排, 第一个点几乎
%  贴着末字, 读起来像多标了一个句点; 页码前也没有余量。两处各补
%  0.5em。
% --------------------------------------------------------------
\makeatletter
\patchcmd{\@dottedtocline}{\nobreak\leaders}{\nobreak\hspace{0.5em}\leaders}%
  {\typeout{MT: TOC 引导点前已留 0.5em}}%
  {\typeout{MT-WARN: TOC 引导点前间距补丁未生效}}
\patchcmd{\@dottedtocline}{\hfill\nobreak\hb@xt@\@pnumwidth}%
  {\hfill\nobreak\hspace{0.5em}\hb@xt@\@pnumwidth}%
  {\typeout{MT: TOC 页码前已留 0.5em}}%
  {\typeout{MT-WARN: TOC 页码前间距补丁未生效}}
\makeatother

% --------------------------------------------------------------
%  正文半角标点后的间距
%  本书体例用半角 , ; : ? ! 与 ASCII 括号, 由 xeCJK 按 HalfRight
%  (右半留白) 处理。放宽到 0.45em, 与全角标点右半留白大致齐平。
% --------------------------------------------------------------
\xeCJKsetup{CJKecglue={\hskip 0.45em plus 0.06em minus 0.09em}}

% --------------------------------------------------------------
%  页眉页脚: 偶数页眉右对齐斜体节题; 页码页脚外侧
% --------------------------------------------------------------
\fancyhf{}
\fancyhead[RE]{\small\itshape\rightmark}
\fancyfoot[LE,RO]{\small\thepage}
\renewcommand{\headrulewidth}{0pt}
\renewcommand{\sectionmark}[1]{\markright{\thesection.\ #1}}
\fancypagestyle{plain}{%
  \fancyhf{}%
  \fancyfoot[LE,RO]{\small \thepage}%
  \renewcommand{\headrulewidth}{0pt}%
}

% --------------------------------------------------------------
%  原书版式记号宏
% --------------------------------------------------------------
% 边码: 原打字稿页码(G. 手稿系统里的自引页码), 排在页边
\usepackage{marginnote}
\newcommand{\origpage}[1]{\marginnote{\normalfont\scriptsize p.~#1}}

% ⋄ 记号: 段首/行内正常大小; 上标形式为作者后补文字的位置标记
\DeclareRobustCommand{\pmark}{\ensuremath{\diamond}\hspace{0.15em}}
\DeclareRobustCommand{\dmark}{\textsuperscript{\ensuremath{\diamond}}}

% 脚注人工编号: \fnote{795}{(*)}{文本} —— 编号与(*)角标均照录原书
\newcommand{\fnote}[3]{%
  \begingroup
  \renewcommand{\thefootnote}{#1#2}%
  \footnote{#3}%
  \endgroup
}

% 第一部分的脚注: 页内序号作上标(#1), 注文行首印 "序号(连续注号)",
% 正文在序号之后另行括注同一连续注号 —— \fnoteP{1}{(3)}{注文}
% 原书有三条注无连续注号而改标日期(ch07 注3 = 30 Septembre,
% ch09 注11/12 = 8 août), 第二参数留空, 日期写进注文之首。
\newcommand{\fnoteP}[3]{%
  \begingroup
  \renewcommand{\thefootnote}{#1}\footnotemark
  \renewcommand{\thefootnote}{#1#2}\footnotetext{#3}%
  \endgroup
}

% 正文呼号与注文行首标号不同(原书照录)
\newcommand{\fnotevariant}[4]{%
  \begingroup
  \renewcommand{\thefootnote}{#1#2}\footnotemark
  \renewcommand{\thefootnote}{#1#3}\footnotetext{#4}%
  \endgroup
}

% 大注块标题: \notehead{171³}{斜体评论} —— 正文随后另起
\newcommand{\notehead}[2]{%
  \par\bigskip
  \noindent{\headfont 注~#1}\ifblank{#2}{}{\hspace{0.5em}{\itshape #2}}%
  \par\nobreak\smallskip\nobreak
}

% ∗ ∗ ∗ 分隔
\newcommand{\asterisms}{%
  \par\vspace{1.2ex}%
  \noindent\hspace*{2em}$\ast$\hspace{1.2em}$\ast$\hspace{1.2em}$\ast$%
  \par\vspace{1.2ex}%
}

% 无编号子目的目录条目
\newcommand{\tocsubsub}[1]{%
  \phantomsection
  \addcontentsline{toc}{subsubsection}{#1}%
}

% 部扉: 仿原书部题页
\newcommand{\partdivider}[3]{%
  \clearpage
  \thispagestyle{empty}%
  \phantomsection
  \addcontentsline{toc}{part}{#1\quad #3}%
  \vspace*{6.5cm}
  \begin{center}
    {\fontsize{16pt}{26pt}\selectfont\headfont #1\par}
    \vspace{0.9cm}
    {\fontsize{22pt}{34pt}\selectfont\headfont #2\par}
  \end{center}
  \clearpage
}

% 献辞页: 仿原书部题后的题献, 独占一页且不入目录
\newcommand{\dedication}[1]{%
  \thispagestyle{empty}%
  \vspace*{4cm}%
  \begin{center}#1\end{center}%
  \clearpage
}

% --------------------------------------------------------------
%  各卷"译者说明"的共用体例段(各卷的范围段自写, 后接本段)
% --------------------------------------------------------------
\newcommand{\CommonNoteBody}{%
为便于与原书及学界通行的打字稿页码对勘, 本册保留:

\begin{itemize}
  \item 原书全部章节编号(章~1--18, 节如~5.2、18.2, 子目如~18.1.1)与原书小节标题;
  \item 原书脚注的编号与角标——
        \begin{itemize}
          \item 前奏(第~1--4 章)与第一部的注在各页内重新编号, 正文中的上标即该
                页序号; 第一部紧随其后的括注(如“(3)”)是该部分的连续注号, 与第~12
                章“第一部分的注”里的“注~3”一一对应, 注文行首照原书印作“1(3)”式。
                其中三条注原书未标连续注号而改标日期(ch7 注~3 = 30 Septembre、
                ch9 注~11/12 = 8 août), 本册照录其日期;
          \item 第二部第~13--16~章各自重新编号, 分别为~1--41、1--86、1--132、1--29;
                第三部第~17~章自~1~起至~12, 第~18~章自~1~重新起算并延续至~794;
                末段自~\textbf{795(*)}~续至~1053(*)。各段与原书一一对应;
        \end{itemize}
  \item 页边码“p.~1”“p.~A3”“p.~1066”式小注——这是作者引用自己打字稿原稿时
        所用的页码, 与本排印本页码不同。前奏四章的边码各有前缀: 第~1 章作
        \texttt{A1…A6}, 第~2 章作 \texttt{P1…P65}, 第~3 章作 \texttt{L1…L56},
        第~4 章作罗马数字 \texttt{I…XXII}, 自第一部起作纯数字 \texttt{1…171};
  \item 作者用以标记段落的~$\diamond$~记号, 以及标示后来补写文字的上标
        ~\textsuperscript{$\diamond$}~记号;
  \item 原书大小注(“注~44'”“注~46$^1$”“注~171³”等)的编号与层级。
\end{itemize}

第~12 章“第一部分的注”把第~5 至第~11 章各页脚注汇为一编, 并给每条注配上节题
(如“12.19. 耶稣与十二使徒”), 因此部分注文与正文脚注重复——这是原书的体例, 非
重出之误。这一部分属随笔式的自省文字, 同一句法语在正文脚注与章末汇编注中
各有措辞, 本册有意保留各自的译笔, 不强求统一。

除上述标记外, 译文力求贴近作者驳杂、跳跃而辛辣的文风; 数学名词的译法见各卷
术语表(如~\emph{topos}~不译, \emph{faisceaux pervers}~作“反常层”)。文中的日期
括注(如“(5~月~30~日)”“(1984年3月补写)”)为作者历次补写时的纪日, 一律照录。文中
另有一些对注号、页码的引用(如“n°~1695”“打字稿第~227~页”)沿袭作者打字稿的编号
系统, 与本排印本不尽一致, 亦无法在本书中落实, 悉照原文录出。

全套五卷(另有全书合订一册)的分工是:

\begin{itemize}
  \item \textbf{卷一·前奏}——主题的呈现: 或四个乐章的前奏(原书 p.23--122);
  \item \textbf{卷二·第一部}——愚妄与更新(原书 p.123--254);
  \item \textbf{卷三·第二部}——葬礼(I): 或中国皇帝的新衣(原书 p.255--412);
  \item \textbf{卷四·第三部}——葬礼(II): 或阴阳之钥(原书 p.413--632);
  \item \textbf{卷五·第四部}——葬礼(III): 或四种运算(原书 p.633--929)。
\end{itemize}
}
                       % 共用导言区
%      % ================================================================
%  《收获与播种》统一封面 —— TikZ 原生矢量绘制(不依赖任何外部图片)
%  母题: 暮色中的麦田 —— 犁沟向地平线的灭点收束, 前景麦穗剪影,
%        一位播种者把一把种子撒向天空, 种子沿弧线落入田里。
%  用法(在 \begin{document} 之前, 依次):
%      \input{common}                       % 共用导言区
%      \input{cover-art}                    % 本文件
%      \CoverSetup{prelude}{卷一}{前奏}%
%                 {主题的呈现——或四个乐章的前奏}{p.23–122}
%  正文中再 \makecover 即可。配色关键字:
%      prelude / part1 / part2 / part3 / part4 / all
%  五卷共用同一幅画, 只换配色与卷次文字, 以保证"一套书"的观感。
% ================================================================

\RequirePackage{tikz}

% --------------------------------------------------------------
%  配色: 六套(前奏 / 四部 / 全书), 各为同一暮色母题的变调
% --------------------------------------------------------------
\newcommand{\CoverPalette}[1]{%
  \ifstrequal{#1}{prelude}{%
    \definecolor{cvSkyTop}{RGB}{10,26,48}%
    \definecolor{cvSkyMid}{RGB}{46,84,116}%
    \definecolor{cvSkyBot}{RGB}{232,164,92}%
    \definecolor{cvSun}{RGB}{255,224,168}%
    \definecolor{cvFieldTop}{RGB}{150,110,52}%
    \definecolor{cvFieldBot}{RGB}{58,36,18}%
    \definecolor{cvFurrow}{RGB}{40,22,10}%
    \definecolor{cvSil}{RGB}{30,18,10}%
    \definecolor{cvCream}{RGB}{247,240,228}%
  }{%
  \ifstrequal{#1}{part1}{%
    \definecolor{cvSkyTop}{RGB}{22,70,100}%
    \definecolor{cvSkyMid}{RGB}{92,150,168}%
    \definecolor{cvSkyBot}{RGB}{240,200,104}%
    \definecolor{cvSun}{RGB}{255,244,200}%
    \definecolor{cvFieldTop}{RGB}{176,140,58}%
    \definecolor{cvFieldBot}{RGB}{84,54,22}%
    \definecolor{cvFurrow}{RGB}{45,28,12}%
    \definecolor{cvSil}{RGB}{34,22,10}%
    \definecolor{cvCream}{RGB}{250,245,232}%
  }{%
  \ifstrequal{#1}{part2}{%
    \definecolor{cvSkyTop}{RGB}{46,20,30}%
    \definecolor{cvSkyMid}{RGB}{128,52,54}%
    \definecolor{cvSkyBot}{RGB}{226,120,64}%
    \definecolor{cvSun}{RGB}{255,206,150}%
    \definecolor{cvFieldTop}{RGB}{132,84,44}%
    \definecolor{cvFieldBot}{RGB}{48,26,18}%
    \definecolor{cvFurrow}{RGB}{36,16,14}%
    \definecolor{cvSil}{RGB}{26,12,12}%
    \definecolor{cvCream}{RGB}{248,236,224}%
  }{%
  \ifstrequal{#1}{part3}{%
    \definecolor{cvSkyTop}{RGB}{30,20,54}%
    \definecolor{cvSkyMid}{RGB}{92,62,124}%
    \definecolor{cvSkyBot}{RGB}{216,132,140}%
    \definecolor{cvSun}{RGB}{255,214,196}%
    \definecolor{cvFieldTop}{RGB}{128,90,74}%
    \definecolor{cvFieldBot}{RGB}{44,26,30}%
    \definecolor{cvFurrow}{RGB}{34,20,30}%
    \definecolor{cvSil}{RGB}{22,14,22}%
    \definecolor{cvCream}{RGB}{246,238,238}%
  }{%
  \ifstrequal{#1}{part4}{%
    \definecolor{cvSkyTop}{RGB}{12,28,28}%
    \definecolor{cvSkyMid}{RGB}{44,86,72}%
    \definecolor{cvSkyBot}{RGB}{186,168,86}%
    \definecolor{cvSun}{RGB}{240,232,168}%
    \definecolor{cvFieldTop}{RGB}{110,104,50}%
    \definecolor{cvFieldBot}{RGB}{30,34,20}%
    \definecolor{cvFurrow}{RGB}{26,32,20}%
    \definecolor{cvSil}{RGB}{14,20,14}%
    \definecolor{cvCream}{RGB}{240,238,226}%
  }{%
    % 默认: 全书(麦浪金褐)
    \definecolor{cvSkyTop}{RGB}{28,40,32}%
    \definecolor{cvSkyMid}{RGB}{128,110,64}%
    \definecolor{cvSkyBot}{RGB}{232,178,88}%
    \definecolor{cvSun}{RGB}{255,236,180}%
    \definecolor{cvFieldTop}{RGB}{168,124,52}%
    \definecolor{cvFieldBot}{RGB}{62,40,20}%
    \definecolor{cvFurrow}{RGB}{42,26,12}%
    \definecolor{cvSil}{RGB}{28,18,10}%
    \definecolor{cvCream}{RGB}{250,242,226}%
  }}}}}%
}

% --------------------------------------------------------------
%  一键设定:  \CoverSetup{配色}{卷次}{卷名}{副题}{原书页码范围}
% --------------------------------------------------------------
\newcommand{\CoverSetup}[5]{%
  \CoverPalette{#1}%
  \newcommand{\VolNum}{#2}%
  \newcommand{\VolName}{#3}%
  \newcommand{\VolSub}{#4}%
  \newcommand{\VolPages}{#5}%
}

% --------------------------------------------------------------
%  一支麦穗:  \cvWheatEar{x}{底 y}{高}{倾角}
%  茎略弯, 麦粒成对左右斜出, 顶端一芒三针
% --------------------------------------------------------------
\newcommand{\cvWheatEar}[4]{%
  \begin{scope}[shift={(#1,#2)}, rotate=#4]
    \draw[cvSil, line width=0.85mm, line cap=round]
      (0,0) .. controls (0.5,{0.35*#3}) and (0.8,{0.72*#3}) .. ({0.55*#3},#3);
    \foreach \cvK in {0,...,6}{%
      \pgfmathsetmacro{\cvT}{0.50+0.50*\cvK/6}
      \pgfmathsetmacro{\cvGY}{\cvT*#3}
      \pgfmathsetmacro{\cvGX}{0.55*\cvT*\cvT*#3}
      \pgfmathsetmacro{\cvGS}{1.0-0.34*\cvT}
      \begin{scope}[shift={(\cvGX,\cvGY)}, rotate=28]
        \fill[cvSil] (0,{1.9*\cvGS}) ellipse [x radius={0.92*\cvGS}, y radius={2.3*\cvGS}];
      \end{scope}
      \begin{scope}[shift={(\cvGX,\cvGY)}, rotate=-28]
        \fill[cvSil] (0,{-1.9*\cvGS}) ellipse [x radius={0.92*\cvGS}, y radius={2.3*\cvGS}];
      \end{scope}
    }
    \begin{scope}[shift={({0.55*#3},#3)}]
      \fill[cvSil] (0,1.5) ellipse [x radius=0.78, y radius=1.95];
    \end{scope}
    \draw[cvSil, line width=0.26mm, line cap=round]
      ({0.55*#3},{#3+2.3}) -- ({0.55*#3-1.7},{#3+9.5});
    \draw[cvSil, line width=0.26mm, line cap=round]
      ({0.55*#3},{#3+2.3}) -- ({0.55*#3+2.1},{#3+10.5});
    \draw[cvSil, line width=0.26mm, line cap=round]
      ({0.55*#3},{#3+2.3}) -- ({0.55*#3},{#3+11.5});
  \end{scope}
}

% --------------------------------------------------------------
%  封面
% --------------------------------------------------------------
\makeatletter
\renewcommand{\makecover}{%
  \begin{titlepage}
    \thispagestyle{empty}
    \begin{tikzpicture}[remember picture, overlay, x=1mm, y=1mm,
                        shift={(current page.south west)}]
      % ===== 天空 =====
      \shade[top color=cvSkyTop, middle color=cvSkyMid, bottom color=cvSkyBot]
        (0,105) rectangle (210,297);
      % ===== 田野 =====
      \shade[top color=cvFieldTop, bottom color=cvFieldBot]
        (0,0) rectangle (210,105);
      % ===== 落日与光晕 =====
      \begin{scope}
        \clip (0,105) rectangle (210,297);
        \fill[cvSun, opacity=0.14] (152,132) circle (46);
        \fill[cvSun, opacity=0.24] (152,132) circle (30);
      \end{scope}
      \fill[cvSun, opacity=0.95] (152,132) circle (16.5);
      % ===== 犁沟: 向地平线灭点收束 =====
      \begin{scope}
        \clip (0,0) rectangle (210,105);
        \foreach \cvI in {1,...,20}{%
          \pgfmathsetmacro{\cvDX}{3.1*pow(\cvI,1.55)}
          \draw[cvFurrow, opacity=0.30, line width=0.32mm] ({105+\cvDX},0) -- (105,105);
          \draw[cvFurrow, opacity=0.30, line width=0.32mm] ({105-\cvDX},0) -- (105,105);
        }
        \foreach \cvY in {8,18,31,47,68,92}{%
          \draw[cvFurrow, opacity=0.15, line width=0.3mm]
            (0,\cvY) .. controls (60,{\cvY+3.5}) and (150,{\cvY-3}) .. (210,{\cvY+1.5});
        }
      \end{scope}
      % 地平线
      \draw[cvSil, opacity=0.45, line width=0.45mm] (0,105) -- (210,105);
      % ===== 播种者剪影(立于地平线) =====
      \begin{scope}[cvSil]
        \draw[line width=1.9mm, line cap=round, line join=round]
          (55.5,116) -- (50,105)  (55.5,116) -- (60.5,105.5);
        \draw[line width=2.3mm, line cap=round, line join=round]
          (55.5,116) -- (55,123.5);
        \draw[line width=1.6mm, line cap=round, line join=round]
          (55.2,123.2) -- (61.5,126.6) -- (68,129.4);
        \draw[line width=1.5mm, line cap=round, line join=round]
          (55.2,123.2) -- (51.4,119.8) -- (48.6,116.2);
        \fill (55.6,127.3) circle (2.35);
      \end{scope}
      % ===== 种子: 沿贝塞尔弧撒出 =====
      \draw[cvCream, opacity=0.13, line width=0.5mm]
        (68,129.4) .. controls (100,168) and (160,162) .. (196,104);
      % 掌心那一把(尚未撒出)
      \foreach \cvI in {1,...,5}{%
        \pgfmathsetmacro{\cvJX}{1.1*sin(\cvI*53)}
        \pgfmathsetmacro{\cvJY}{1.1*cos(\cvI*83)}
        \fill[cvCream, opacity=0.9] ({68+\cvJX},{129.4+\cvJY}) circle (1.0);
      }
      \foreach \cvI in {1,...,26}{%
        \pgfmathsetmacro{\cvT}{0.035+0.92*\cvI/26}
        \pgfmathsetmacro{\cvU}{1-\cvT}
        \pgfmathsetmacro{\cvSX}{\cvU*\cvU*\cvU*68 + 3*\cvU*\cvU*\cvT*100
                                + 3*\cvU*\cvT*\cvT*160 + \cvT*\cvT*\cvT*196}
        \pgfmathsetmacro{\cvSY}{\cvU*\cvU*\cvU*129.4 + 3*\cvU*\cvU*\cvT*168
                                + 3*\cvU*\cvT*\cvT*162 + \cvT*\cvT*\cvT*104}
        \pgfmathsetmacro{\cvR}{1.55-0.85*\cvT}
        \pgfmathsetmacro{\cvJX}{1.7*sin(\cvI*97)}
        \pgfmathsetmacro{\cvJY}{1.7*cos(\cvI*71)}
        \fill[cvCream, opacity={0.88-0.28*\cvT}] ({\cvSX+\cvJX},{\cvSY+\cvJY}) circle (\cvR);
      }
      % ===== 前景麦穗(两排) =====
      \foreach \cvI in {1,...,11}{%
        \pgfmathsetmacro{\cvWX}{4+19*(\cvI-1)+1.6*sin(\cvI*2.1)}
        \pgfmathsetmacro{\cvWH}{23+7*sin(\cvI*1.7)}
        \pgfmathsetmacro{\cvWL}{3.2*sin(\cvI*2.9)}
        \begin{scope}[opacity=0.45]
          \cvWheatEar{\cvWX}{30}{\cvWH}{\cvWL}
        \end{scope}
      }
      \foreach \cvI in {1,...,10}{%
        \pgfmathsetmacro{\cvWX}{9+20.4*(\cvI-1)+1.5*sin(\cvI*1.3)}
        \pgfmathsetmacro{\cvWH}{41+10*sin(\cvI*2.3)}
        \pgfmathsetmacro{\cvWL}{4.5*sin(\cvI*3.7)}
        \cvWheatEar{\cvWX}{30}{\cvWH}{\cvWL}
      }
      % ===== 压暗一层, 统一色调并保证文字可读 =====
      \fill[black, opacity=0.10] (0,105) rectangle (210,297);
      % ===== 底部信息带 =====
      \fill[cvSil, opacity=0.72] (0,0) rectangle (210,30);
      \draw[cvCream, opacity=0.40, line width=0.4mm] (0,30) -- (210,30);
      % ===== 文字 =====
      \node[anchor=north west, inner sep=0, text=cvCream] at (22,268)
        {\fontsize{38pt}{45pt}\selectfont\headfont \ziju{0.08}\BookTitleCN};
      \node[anchor=north west, inner sep=0, text=cvCream] at (22.6,242)
        {\fontsize{15pt}{20pt}\selectfont\itshape \BookTitleEN};
      \draw[cvCream, opacity=0.55, line width=0.5mm] (22,233) -- (120,233);
      \node[anchor=north west, inner sep=0, text=cvCream] at (22,224)
        {\fontsize{11pt}{15pt}\selectfont 一位数学家对往昔的反思与见证};
      \node[anchor=north west, inner sep=0, text=cvCream, opacity=0.78] at (22.4,208)
        {\fontsize{10pt}{14pt}\selectfont \ziju{0.14}\VolNum};
      \node[anchor=north west, inner sep=0, text=cvCream] at (22,199)
        {\fontsize{30pt}{37pt}\selectfont\headfont \ziju{0.1}\VolName};
      \node[anchor=north west, inner sep=0, text=cvCream, opacity=0.92] at (22,180)
        {\fontsize{13pt}{18pt}\selectfont \VolSub};
      \node[anchor=north west, inner sep=0, text=cvCream] at (22,19)
        {\fontsize{11.5pt}{15pt}\selectfont \OriginalAuthor\quad 著};
      \node[anchor=north west, inner sep=0, text=cvCream, opacity=0.85] at (22,7)
        {\fontsize{9.5pt}{13pt}\selectfont \Translator\ 译\quad ·\quad
         \ifdefempty{\TranslationDate}{\today}{\TranslationDate}};
      \node[anchor=north east, inner sep=0, text=cvCream, opacity=0.85] at (188,13)
        {\fontsize{10pt}{14pt}\selectfont 原书 \VolPages};
    \end{tikzpicture}
  \end{titlepage}%
}
\makeatother
                    % 本文件
%      \CoverSetup{prelude}{卷一}{前奏}%
%                 {主题的呈现——或四个乐章的前奏}{p.23–122}
%  正文中再 \makecover 即可。配色关键字:
%      prelude / part1 / part2 / part3 / part4 / all
%  五卷共用同一幅画, 只换配色与卷次文字, 以保证"一套书"的观感。
% ================================================================

\RequirePackage{tikz}

% --------------------------------------------------------------
%  配色: 六套(前奏 / 四部 / 全书), 各为同一暮色母题的变调
% --------------------------------------------------------------
\newcommand{\CoverPalette}[1]{%
  \ifstrequal{#1}{prelude}{%
    \definecolor{cvSkyTop}{RGB}{10,26,48}%
    \definecolor{cvSkyMid}{RGB}{46,84,116}%
    \definecolor{cvSkyBot}{RGB}{232,164,92}%
    \definecolor{cvSun}{RGB}{255,224,168}%
    \definecolor{cvFieldTop}{RGB}{150,110,52}%
    \definecolor{cvFieldBot}{RGB}{58,36,18}%
    \definecolor{cvFurrow}{RGB}{40,22,10}%
    \definecolor{cvSil}{RGB}{30,18,10}%
    \definecolor{cvCream}{RGB}{247,240,228}%
  }{%
  \ifstrequal{#1}{part1}{%
    \definecolor{cvSkyTop}{RGB}{22,70,100}%
    \definecolor{cvSkyMid}{RGB}{92,150,168}%
    \definecolor{cvSkyBot}{RGB}{240,200,104}%
    \definecolor{cvSun}{RGB}{255,244,200}%
    \definecolor{cvFieldTop}{RGB}{176,140,58}%
    \definecolor{cvFieldBot}{RGB}{84,54,22}%
    \definecolor{cvFurrow}{RGB}{45,28,12}%
    \definecolor{cvSil}{RGB}{34,22,10}%
    \definecolor{cvCream}{RGB}{250,245,232}%
  }{%
  \ifstrequal{#1}{part2}{%
    \definecolor{cvSkyTop}{RGB}{46,20,30}%
    \definecolor{cvSkyMid}{RGB}{128,52,54}%
    \definecolor{cvSkyBot}{RGB}{226,120,64}%
    \definecolor{cvSun}{RGB}{255,206,150}%
    \definecolor{cvFieldTop}{RGB}{132,84,44}%
    \definecolor{cvFieldBot}{RGB}{48,26,18}%
    \definecolor{cvFurrow}{RGB}{36,16,14}%
    \definecolor{cvSil}{RGB}{26,12,12}%
    \definecolor{cvCream}{RGB}{248,236,224}%
  }{%
  \ifstrequal{#1}{part3}{%
    \definecolor{cvSkyTop}{RGB}{30,20,54}%
    \definecolor{cvSkyMid}{RGB}{92,62,124}%
    \definecolor{cvSkyBot}{RGB}{216,132,140}%
    \definecolor{cvSun}{RGB}{255,214,196}%
    \definecolor{cvFieldTop}{RGB}{128,90,74}%
    \definecolor{cvFieldBot}{RGB}{44,26,30}%
    \definecolor{cvFurrow}{RGB}{34,20,30}%
    \definecolor{cvSil}{RGB}{22,14,22}%
    \definecolor{cvCream}{RGB}{246,238,238}%
  }{%
  \ifstrequal{#1}{part4}{%
    \definecolor{cvSkyTop}{RGB}{12,28,28}%
    \definecolor{cvSkyMid}{RGB}{44,86,72}%
    \definecolor{cvSkyBot}{RGB}{186,168,86}%
    \definecolor{cvSun}{RGB}{240,232,168}%
    \definecolor{cvFieldTop}{RGB}{110,104,50}%
    \definecolor{cvFieldBot}{RGB}{30,34,20}%
    \definecolor{cvFurrow}{RGB}{26,32,20}%
    \definecolor{cvSil}{RGB}{14,20,14}%
    \definecolor{cvCream}{RGB}{240,238,226}%
  }{%
    % 默认: 全书(麦浪金褐)
    \definecolor{cvSkyTop}{RGB}{28,40,32}%
    \definecolor{cvSkyMid}{RGB}{128,110,64}%
    \definecolor{cvSkyBot}{RGB}{232,178,88}%
    \definecolor{cvSun}{RGB}{255,236,180}%
    \definecolor{cvFieldTop}{RGB}{168,124,52}%
    \definecolor{cvFieldBot}{RGB}{62,40,20}%
    \definecolor{cvFurrow}{RGB}{42,26,12}%
    \definecolor{cvSil}{RGB}{28,18,10}%
    \definecolor{cvCream}{RGB}{250,242,226}%
  }}}}}%
}

% --------------------------------------------------------------
%  一键设定:  \CoverSetup{配色}{卷次}{卷名}{副题}{原书页码范围}
% --------------------------------------------------------------
\newcommand{\CoverSetup}[5]{%
  \CoverPalette{#1}%
  \newcommand{\VolNum}{#2}%
  \newcommand{\VolName}{#3}%
  \newcommand{\VolSub}{#4}%
  \newcommand{\VolPages}{#5}%
}

% --------------------------------------------------------------
%  一支麦穗:  \cvWheatEar{x}{底 y}{高}{倾角}
%  茎略弯, 麦粒成对左右斜出, 顶端一芒三针
% --------------------------------------------------------------
\newcommand{\cvWheatEar}[4]{%
  \begin{scope}[shift={(#1,#2)}, rotate=#4]
    \draw[cvSil, line width=0.85mm, line cap=round]
      (0,0) .. controls (0.5,{0.35*#3}) and (0.8,{0.72*#3}) .. ({0.55*#3},#3);
    \foreach \cvK in {0,...,6}{%
      \pgfmathsetmacro{\cvT}{0.50+0.50*\cvK/6}
      \pgfmathsetmacro{\cvGY}{\cvT*#3}
      \pgfmathsetmacro{\cvGX}{0.55*\cvT*\cvT*#3}
      \pgfmathsetmacro{\cvGS}{1.0-0.34*\cvT}
      \begin{scope}[shift={(\cvGX,\cvGY)}, rotate=28]
        \fill[cvSil] (0,{1.9*\cvGS}) ellipse [x radius={0.92*\cvGS}, y radius={2.3*\cvGS}];
      \end{scope}
      \begin{scope}[shift={(\cvGX,\cvGY)}, rotate=-28]
        \fill[cvSil] (0,{-1.9*\cvGS}) ellipse [x radius={0.92*\cvGS}, y radius={2.3*\cvGS}];
      \end{scope}
    }
    \begin{scope}[shift={({0.55*#3},#3)}]
      \fill[cvSil] (0,1.5) ellipse [x radius=0.78, y radius=1.95];
    \end{scope}
    \draw[cvSil, line width=0.26mm, line cap=round]
      ({0.55*#3},{#3+2.3}) -- ({0.55*#3-1.7},{#3+9.5});
    \draw[cvSil, line width=0.26mm, line cap=round]
      ({0.55*#3},{#3+2.3}) -- ({0.55*#3+2.1},{#3+10.5});
    \draw[cvSil, line width=0.26mm, line cap=round]
      ({0.55*#3},{#3+2.3}) -- ({0.55*#3},{#3+11.5});
  \end{scope}
}

% --------------------------------------------------------------
%  封面
% --------------------------------------------------------------
\makeatletter
\renewcommand{\makecover}{%
  \begin{titlepage}
    \thispagestyle{empty}
    \begin{tikzpicture}[remember picture, overlay, x=1mm, y=1mm,
                        shift={(current page.south west)}]
      % ===== 天空 =====
      \shade[top color=cvSkyTop, middle color=cvSkyMid, bottom color=cvSkyBot]
        (0,105) rectangle (210,297);
      % ===== 田野 =====
      \shade[top color=cvFieldTop, bottom color=cvFieldBot]
        (0,0) rectangle (210,105);
      % ===== 落日与光晕 =====
      \begin{scope}
        \clip (0,105) rectangle (210,297);
        \fill[cvSun, opacity=0.14] (152,132) circle (46);
        \fill[cvSun, opacity=0.24] (152,132) circle (30);
      \end{scope}
      \fill[cvSun, opacity=0.95] (152,132) circle (16.5);
      % ===== 犁沟: 向地平线灭点收束 =====
      \begin{scope}
        \clip (0,0) rectangle (210,105);
        \foreach \cvI in {1,...,20}{%
          \pgfmathsetmacro{\cvDX}{3.1*pow(\cvI,1.55)}
          \draw[cvFurrow, opacity=0.30, line width=0.32mm] ({105+\cvDX},0) -- (105,105);
          \draw[cvFurrow, opacity=0.30, line width=0.32mm] ({105-\cvDX},0) -- (105,105);
        }
        \foreach \cvY in {8,18,31,47,68,92}{%
          \draw[cvFurrow, opacity=0.15, line width=0.3mm]
            (0,\cvY) .. controls (60,{\cvY+3.5}) and (150,{\cvY-3}) .. (210,{\cvY+1.5});
        }
      \end{scope}
      % 地平线
      \draw[cvSil, opacity=0.45, line width=0.45mm] (0,105) -- (210,105);
      % ===== 播种者剪影(立于地平线) =====
      \begin{scope}[cvSil]
        \draw[line width=1.9mm, line cap=round, line join=round]
          (55.5,116) -- (50,105)  (55.5,116) -- (60.5,105.5);
        \draw[line width=2.3mm, line cap=round, line join=round]
          (55.5,116) -- (55,123.5);
        \draw[line width=1.6mm, line cap=round, line join=round]
          (55.2,123.2) -- (61.5,126.6) -- (68,129.4);
        \draw[line width=1.5mm, line cap=round, line join=round]
          (55.2,123.2) -- (51.4,119.8) -- (48.6,116.2);
        \fill (55.6,127.3) circle (2.35);
      \end{scope}
      % ===== 种子: 沿贝塞尔弧撒出 =====
      \draw[cvCream, opacity=0.13, line width=0.5mm]
        (68,129.4) .. controls (100,168) and (160,162) .. (196,104);
      % 掌心那一把(尚未撒出)
      \foreach \cvI in {1,...,5}{%
        \pgfmathsetmacro{\cvJX}{1.1*sin(\cvI*53)}
        \pgfmathsetmacro{\cvJY}{1.1*cos(\cvI*83)}
        \fill[cvCream, opacity=0.9] ({68+\cvJX},{129.4+\cvJY}) circle (1.0);
      }
      \foreach \cvI in {1,...,26}{%
        \pgfmathsetmacro{\cvT}{0.035+0.92*\cvI/26}
        \pgfmathsetmacro{\cvU}{1-\cvT}
        \pgfmathsetmacro{\cvSX}{\cvU*\cvU*\cvU*68 + 3*\cvU*\cvU*\cvT*100
                                + 3*\cvU*\cvT*\cvT*160 + \cvT*\cvT*\cvT*196}
        \pgfmathsetmacro{\cvSY}{\cvU*\cvU*\cvU*129.4 + 3*\cvU*\cvU*\cvT*168
                                + 3*\cvU*\cvT*\cvT*162 + \cvT*\cvT*\cvT*104}
        \pgfmathsetmacro{\cvR}{1.55-0.85*\cvT}
        \pgfmathsetmacro{\cvJX}{1.7*sin(\cvI*97)}
        \pgfmathsetmacro{\cvJY}{1.7*cos(\cvI*71)}
        \fill[cvCream, opacity={0.88-0.28*\cvT}] ({\cvSX+\cvJX},{\cvSY+\cvJY}) circle (\cvR);
      }
      % ===== 前景麦穗(两排) =====
      \foreach \cvI in {1,...,11}{%
        \pgfmathsetmacro{\cvWX}{4+19*(\cvI-1)+1.6*sin(\cvI*2.1)}
        \pgfmathsetmacro{\cvWH}{23+7*sin(\cvI*1.7)}
        \pgfmathsetmacro{\cvWL}{3.2*sin(\cvI*2.9)}
        \begin{scope}[opacity=0.45]
          \cvWheatEar{\cvWX}{30}{\cvWH}{\cvWL}
        \end{scope}
      }
      \foreach \cvI in {1,...,10}{%
        \pgfmathsetmacro{\cvWX}{9+20.4*(\cvI-1)+1.5*sin(\cvI*1.3)}
        \pgfmathsetmacro{\cvWH}{41+10*sin(\cvI*2.3)}
        \pgfmathsetmacro{\cvWL}{4.5*sin(\cvI*3.7)}
        \cvWheatEar{\cvWX}{30}{\cvWH}{\cvWL}
      }
      % ===== 压暗一层, 统一色调并保证文字可读 =====
      \fill[black, opacity=0.10] (0,105) rectangle (210,297);
      % ===== 底部信息带 =====
      \fill[cvSil, opacity=0.72] (0,0) rectangle (210,30);
      \draw[cvCream, opacity=0.40, line width=0.4mm] (0,30) -- (210,30);
      % ===== 文字 =====
      \node[anchor=north west, inner sep=0, text=cvCream] at (22,268)
        {\fontsize{38pt}{45pt}\selectfont\headfont \ziju{0.08}\BookTitleCN};
      \node[anchor=north west, inner sep=0, text=cvCream] at (22.6,242)
        {\fontsize{15pt}{20pt}\selectfont\itshape \BookTitleEN};
      \draw[cvCream, opacity=0.55, line width=0.5mm] (22,233) -- (120,233);
      \node[anchor=north west, inner sep=0, text=cvCream] at (22,224)
        {\fontsize{11pt}{15pt}\selectfont 一位数学家对往昔的反思与见证};
      \node[anchor=north west, inner sep=0, text=cvCream, opacity=0.78] at (22.4,208)
        {\fontsize{10pt}{14pt}\selectfont \ziju{0.14}\VolNum};
      \node[anchor=north west, inner sep=0, text=cvCream] at (22,199)
        {\fontsize{30pt}{37pt}\selectfont\headfont \ziju{0.1}\VolName};
      \node[anchor=north west, inner sep=0, text=cvCream, opacity=0.92] at (22,180)
        {\fontsize{13pt}{18pt}\selectfont \VolSub};
      \node[anchor=north west, inner sep=0, text=cvCream] at (22,19)
        {\fontsize{11.5pt}{15pt}\selectfont \OriginalAuthor\quad 著};
      \node[anchor=north west, inner sep=0, text=cvCream, opacity=0.85] at (22,7)
        {\fontsize{9.5pt}{13pt}\selectfont \Translator\ 译\quad ·\quad
         \ifdefempty{\TranslationDate}{\today}{\TranslationDate}};
      \node[anchor=north east, inner sep=0, text=cvCream, opacity=0.85] at (188,13)
        {\fontsize{10pt}{14pt}\selectfont 原书 \VolPages};
    \end{tikzpicture}
  \end{titlepage}%
}
\makeatother
                    % 本文件
%      \CoverSetup{prelude}{卷一}{前奏}%
%                 {主题的呈现——或四个乐章的前奏}{p.23–122}
%  正文中再 \makecover 即可。配色关键字:
%      prelude / part1 / part2 / part3 / part4 / all
%  五卷共用同一幅画, 只换配色与卷次文字, 以保证"一套书"的观感。
% ================================================================

\RequirePackage{tikz}

% --------------------------------------------------------------
%  配色: 六套(前奏 / 四部 / 全书), 各为同一暮色母题的变调
% --------------------------------------------------------------
\newcommand{\CoverPalette}[1]{%
  \ifstrequal{#1}{prelude}{%
    \definecolor{cvSkyTop}{RGB}{10,26,48}%
    \definecolor{cvSkyMid}{RGB}{46,84,116}%
    \definecolor{cvSkyBot}{RGB}{232,164,92}%
    \definecolor{cvSun}{RGB}{255,224,168}%
    \definecolor{cvFieldTop}{RGB}{150,110,52}%
    \definecolor{cvFieldBot}{RGB}{58,36,18}%
    \definecolor{cvFurrow}{RGB}{40,22,10}%
    \definecolor{cvSil}{RGB}{30,18,10}%
    \definecolor{cvCream}{RGB}{247,240,228}%
  }{%
  \ifstrequal{#1}{part1}{%
    \definecolor{cvSkyTop}{RGB}{22,70,100}%
    \definecolor{cvSkyMid}{RGB}{92,150,168}%
    \definecolor{cvSkyBot}{RGB}{240,200,104}%
    \definecolor{cvSun}{RGB}{255,244,200}%
    \definecolor{cvFieldTop}{RGB}{176,140,58}%
    \definecolor{cvFieldBot}{RGB}{84,54,22}%
    \definecolor{cvFurrow}{RGB}{45,28,12}%
    \definecolor{cvSil}{RGB}{34,22,10}%
    \definecolor{cvCream}{RGB}{250,245,232}%
  }{%
  \ifstrequal{#1}{part2}{%
    \definecolor{cvSkyTop}{RGB}{46,20,30}%
    \definecolor{cvSkyMid}{RGB}{128,52,54}%
    \definecolor{cvSkyBot}{RGB}{226,120,64}%
    \definecolor{cvSun}{RGB}{255,206,150}%
    \definecolor{cvFieldTop}{RGB}{132,84,44}%
    \definecolor{cvFieldBot}{RGB}{48,26,18}%
    \definecolor{cvFurrow}{RGB}{36,16,14}%
    \definecolor{cvSil}{RGB}{26,12,12}%
    \definecolor{cvCream}{RGB}{248,236,224}%
  }{%
  \ifstrequal{#1}{part3}{%
    \definecolor{cvSkyTop}{RGB}{30,20,54}%
    \definecolor{cvSkyMid}{RGB}{92,62,124}%
    \definecolor{cvSkyBot}{RGB}{216,132,140}%
    \definecolor{cvSun}{RGB}{255,214,196}%
    \definecolor{cvFieldTop}{RGB}{128,90,74}%
    \definecolor{cvFieldBot}{RGB}{44,26,30}%
    \definecolor{cvFurrow}{RGB}{34,20,30}%
    \definecolor{cvSil}{RGB}{22,14,22}%
    \definecolor{cvCream}{RGB}{246,238,238}%
  }{%
  \ifstrequal{#1}{part4}{%
    \definecolor{cvSkyTop}{RGB}{12,28,28}%
    \definecolor{cvSkyMid}{RGB}{44,86,72}%
    \definecolor{cvSkyBot}{RGB}{186,168,86}%
    \definecolor{cvSun}{RGB}{240,232,168}%
    \definecolor{cvFieldTop}{RGB}{110,104,50}%
    \definecolor{cvFieldBot}{RGB}{30,34,20}%
    \definecolor{cvFurrow}{RGB}{26,32,20}%
    \definecolor{cvSil}{RGB}{14,20,14}%
    \definecolor{cvCream}{RGB}{240,238,226}%
  }{%
    % 默认: 全书(麦浪金褐)
    \definecolor{cvSkyTop}{RGB}{28,40,32}%
    \definecolor{cvSkyMid}{RGB}{128,110,64}%
    \definecolor{cvSkyBot}{RGB}{232,178,88}%
    \definecolor{cvSun}{RGB}{255,236,180}%
    \definecolor{cvFieldTop}{RGB}{168,124,52}%
    \definecolor{cvFieldBot}{RGB}{62,40,20}%
    \definecolor{cvFurrow}{RGB}{42,26,12}%
    \definecolor{cvSil}{RGB}{28,18,10}%
    \definecolor{cvCream}{RGB}{250,242,226}%
  }}}}}%
}

% --------------------------------------------------------------
%  一键设定:  \CoverSetup{配色}{卷次}{卷名}{副题}{原书页码范围}
% --------------------------------------------------------------
\newcommand{\CoverSetup}[5]{%
  \CoverPalette{#1}%
  \newcommand{\VolNum}{#2}%
  \newcommand{\VolName}{#3}%
  \newcommand{\VolSub}{#4}%
  \newcommand{\VolPages}{#5}%
}

% --------------------------------------------------------------
%  一支麦穗:  \cvWheatEar{x}{底 y}{高}{倾角}
%  茎略弯, 麦粒成对左右斜出, 顶端一芒三针
% --------------------------------------------------------------
\newcommand{\cvWheatEar}[4]{%
  \begin{scope}[shift={(#1,#2)}, rotate=#4]
    \draw[cvSil, line width=0.85mm, line cap=round]
      (0,0) .. controls (0.5,{0.35*#3}) and (0.8,{0.72*#3}) .. ({0.55*#3},#3);
    \foreach \cvK in {0,...,6}{%
      \pgfmathsetmacro{\cvT}{0.50+0.50*\cvK/6}
      \pgfmathsetmacro{\cvGY}{\cvT*#3}
      \pgfmathsetmacro{\cvGX}{0.55*\cvT*\cvT*#3}
      \pgfmathsetmacro{\cvGS}{1.0-0.34*\cvT}
      \begin{scope}[shift={(\cvGX,\cvGY)}, rotate=28]
        \fill[cvSil] (0,{1.9*\cvGS}) ellipse [x radius={0.92*\cvGS}, y radius={2.3*\cvGS}];
      \end{scope}
      \begin{scope}[shift={(\cvGX,\cvGY)}, rotate=-28]
        \fill[cvSil] (0,{-1.9*\cvGS}) ellipse [x radius={0.92*\cvGS}, y radius={2.3*\cvGS}];
      \end{scope}
    }
    \begin{scope}[shift={({0.55*#3},#3)}]
      \fill[cvSil] (0,1.5) ellipse [x radius=0.78, y radius=1.95];
    \end{scope}
    \draw[cvSil, line width=0.26mm, line cap=round]
      ({0.55*#3},{#3+2.3}) -- ({0.55*#3-1.7},{#3+9.5});
    \draw[cvSil, line width=0.26mm, line cap=round]
      ({0.55*#3},{#3+2.3}) -- ({0.55*#3+2.1},{#3+10.5});
    \draw[cvSil, line width=0.26mm, line cap=round]
      ({0.55*#3},{#3+2.3}) -- ({0.55*#3},{#3+11.5});
  \end{scope}
}

% --------------------------------------------------------------
%  封面
% --------------------------------------------------------------
\makeatletter
\renewcommand{\makecover}{%
  \begin{titlepage}
    \thispagestyle{empty}
    \begin{tikzpicture}[remember picture, overlay, x=1mm, y=1mm,
                        shift={(current page.south west)}]
      % ===== 天空 =====
      \shade[top color=cvSkyTop, middle color=cvSkyMid, bottom color=cvSkyBot]
        (0,105) rectangle (210,297);
      % ===== 田野 =====
      \shade[top color=cvFieldTop, bottom color=cvFieldBot]
        (0,0) rectangle (210,105);
      % ===== 落日与光晕 =====
      \begin{scope}
        \clip (0,105) rectangle (210,297);
        \fill[cvSun, opacity=0.14] (152,132) circle (46);
        \fill[cvSun, opacity=0.24] (152,132) circle (30);
      \end{scope}
      \fill[cvSun, opacity=0.95] (152,132) circle (16.5);
      % ===== 犁沟: 向地平线灭点收束 =====
      \begin{scope}
        \clip (0,0) rectangle (210,105);
        \foreach \cvI in {1,...,20}{%
          \pgfmathsetmacro{\cvDX}{3.1*pow(\cvI,1.55)}
          \draw[cvFurrow, opacity=0.30, line width=0.32mm] ({105+\cvDX},0) -- (105,105);
          \draw[cvFurrow, opacity=0.30, line width=0.32mm] ({105-\cvDX},0) -- (105,105);
        }
        \foreach \cvY in {8,18,31,47,68,92}{%
          \draw[cvFurrow, opacity=0.15, line width=0.3mm]
            (0,\cvY) .. controls (60,{\cvY+3.5}) and (150,{\cvY-3}) .. (210,{\cvY+1.5});
        }
      \end{scope}
      % 地平线
      \draw[cvSil, opacity=0.45, line width=0.45mm] (0,105) -- (210,105);
      % ===== 播种者剪影(立于地平线) =====
      \begin{scope}[cvSil]
        \draw[line width=1.9mm, line cap=round, line join=round]
          (55.5,116) -- (50,105)  (55.5,116) -- (60.5,105.5);
        \draw[line width=2.3mm, line cap=round, line join=round]
          (55.5,116) -- (55,123.5);
        \draw[line width=1.6mm, line cap=round, line join=round]
          (55.2,123.2) -- (61.5,126.6) -- (68,129.4);
        \draw[line width=1.5mm, line cap=round, line join=round]
          (55.2,123.2) -- (51.4,119.8) -- (48.6,116.2);
        \fill (55.6,127.3) circle (2.35);
      \end{scope}
      % ===== 种子: 沿贝塞尔弧撒出 =====
      \draw[cvCream, opacity=0.13, line width=0.5mm]
        (68,129.4) .. controls (100,168) and (160,162) .. (196,104);
      % 掌心那一把(尚未撒出)
      \foreach \cvI in {1,...,5}{%
        \pgfmathsetmacro{\cvJX}{1.1*sin(\cvI*53)}
        \pgfmathsetmacro{\cvJY}{1.1*cos(\cvI*83)}
        \fill[cvCream, opacity=0.9] ({68+\cvJX},{129.4+\cvJY}) circle (1.0);
      }
      \foreach \cvI in {1,...,26}{%
        \pgfmathsetmacro{\cvT}{0.035+0.92*\cvI/26}
        \pgfmathsetmacro{\cvU}{1-\cvT}
        \pgfmathsetmacro{\cvSX}{\cvU*\cvU*\cvU*68 + 3*\cvU*\cvU*\cvT*100
                                + 3*\cvU*\cvT*\cvT*160 + \cvT*\cvT*\cvT*196}
        \pgfmathsetmacro{\cvSY}{\cvU*\cvU*\cvU*129.4 + 3*\cvU*\cvU*\cvT*168
                                + 3*\cvU*\cvT*\cvT*162 + \cvT*\cvT*\cvT*104}
        \pgfmathsetmacro{\cvR}{1.55-0.85*\cvT}
        \pgfmathsetmacro{\cvJX}{1.7*sin(\cvI*97)}
        \pgfmathsetmacro{\cvJY}{1.7*cos(\cvI*71)}
        \fill[cvCream, opacity={0.88-0.28*\cvT}] ({\cvSX+\cvJX},{\cvSY+\cvJY}) circle (\cvR);
      }
      % ===== 前景麦穗(两排) =====
      \foreach \cvI in {1,...,11}{%
        \pgfmathsetmacro{\cvWX}{4+19*(\cvI-1)+1.6*sin(\cvI*2.1)}
        \pgfmathsetmacro{\cvWH}{23+7*sin(\cvI*1.7)}
        \pgfmathsetmacro{\cvWL}{3.2*sin(\cvI*2.9)}
        \begin{scope}[opacity=0.45]
          \cvWheatEar{\cvWX}{30}{\cvWH}{\cvWL}
        \end{scope}
      }
      \foreach \cvI in {1,...,10}{%
        \pgfmathsetmacro{\cvWX}{9+20.4*(\cvI-1)+1.5*sin(\cvI*1.3)}
        \pgfmathsetmacro{\cvWH}{41+10*sin(\cvI*2.3)}
        \pgfmathsetmacro{\cvWL}{4.5*sin(\cvI*3.7)}
        \cvWheatEar{\cvWX}{30}{\cvWH}{\cvWL}
      }
      % ===== 压暗一层, 统一色调并保证文字可读 =====
      \fill[black, opacity=0.10] (0,105) rectangle (210,297);
      % ===== 底部信息带 =====
      \fill[cvSil, opacity=0.72] (0,0) rectangle (210,30);
      \draw[cvCream, opacity=0.40, line width=0.4mm] (0,30) -- (210,30);
      % ===== 文字 =====
      \node[anchor=north west, inner sep=0, text=cvCream] at (22,268)
        {\fontsize{38pt}{45pt}\selectfont\headfont \ziju{0.08}\BookTitleCN};
      \node[anchor=north west, inner sep=0, text=cvCream] at (22.6,242)
        {\fontsize{15pt}{20pt}\selectfont\itshape \BookTitleEN};
      \draw[cvCream, opacity=0.55, line width=0.5mm] (22,233) -- (120,233);
      \node[anchor=north west, inner sep=0, text=cvCream] at (22,224)
        {\fontsize{11pt}{15pt}\selectfont 一位数学家对往昔的反思与见证};
      \node[anchor=north west, inner sep=0, text=cvCream, opacity=0.78] at (22.4,208)
        {\fontsize{10pt}{14pt}\selectfont \ziju{0.14}\VolNum};
      \node[anchor=north west, inner sep=0, text=cvCream] at (22,199)
        {\fontsize{30pt}{37pt}\selectfont\headfont \ziju{0.1}\VolName};
      \node[anchor=north west, inner sep=0, text=cvCream, opacity=0.92] at (22,180)
        {\fontsize{13pt}{18pt}\selectfont \VolSub};
      \node[anchor=north west, inner sep=0, text=cvCream] at (22,19)
        {\fontsize{11.5pt}{15pt}\selectfont \OriginalAuthor\quad 著};
      \node[anchor=north west, inner sep=0, text=cvCream, opacity=0.85] at (22,7)
        {\fontsize{9.5pt}{13pt}\selectfont \Translator\ 译\quad ·\quad
         \ifdefempty{\TranslationDate}{\today}{\TranslationDate}};
      \node[anchor=north east, inner sep=0, text=cvCream, opacity=0.85] at (188,13)
        {\fontsize{10pt}{14pt}\selectfont 原书 \VolPages};
    \end{tikzpicture}
  \end{titlepage}%
}
\makeatother
                    % 本文件
%      \CoverSetup{prelude}{卷一}{前奏}%
%                 {主题的呈现——或四个乐章的前奏}{p.23–122}
%  正文中再 \makecover 即可。配色关键字:
%      prelude / part1 / part2 / part3 / part4 / all
%  五卷共用同一幅画, 只换配色与卷次文字, 以保证"一套书"的观感。
% ================================================================

\RequirePackage{tikz}

% --------------------------------------------------------------
%  配色: 六套(前奏 / 四部 / 全书), 各为同一暮色母题的变调
% --------------------------------------------------------------
\newcommand{\CoverPalette}[1]{%
  \ifstrequal{#1}{prelude}{%
    \definecolor{cvSkyTop}{RGB}{10,26,48}%
    \definecolor{cvSkyMid}{RGB}{46,84,116}%
    \definecolor{cvSkyBot}{RGB}{232,164,92}%
    \definecolor{cvSun}{RGB}{255,224,168}%
    \definecolor{cvFieldTop}{RGB}{150,110,52}%
    \definecolor{cvFieldBot}{RGB}{58,36,18}%
    \definecolor{cvFurrow}{RGB}{40,22,10}%
    \definecolor{cvSil}{RGB}{30,18,10}%
    \definecolor{cvCream}{RGB}{247,240,228}%
  }{%
  \ifstrequal{#1}{part1}{%
    \definecolor{cvSkyTop}{RGB}{22,70,100}%
    \definecolor{cvSkyMid}{RGB}{92,150,168}%
    \definecolor{cvSkyBot}{RGB}{240,200,104}%
    \definecolor{cvSun}{RGB}{255,244,200}%
    \definecolor{cvFieldTop}{RGB}{176,140,58}%
    \definecolor{cvFieldBot}{RGB}{84,54,22}%
    \definecolor{cvFurrow}{RGB}{45,28,12}%
    \definecolor{cvSil}{RGB}{34,22,10}%
    \definecolor{cvCream}{RGB}{250,245,232}%
  }{%
  \ifstrequal{#1}{part2}{%
    \definecolor{cvSkyTop}{RGB}{46,20,30}%
    \definecolor{cvSkyMid}{RGB}{128,52,54}%
    \definecolor{cvSkyBot}{RGB}{226,120,64}%
    \definecolor{cvSun}{RGB}{255,206,150}%
    \definecolor{cvFieldTop}{RGB}{132,84,44}%
    \definecolor{cvFieldBot}{RGB}{48,26,18}%
    \definecolor{cvFurrow}{RGB}{36,16,14}%
    \definecolor{cvSil}{RGB}{26,12,12}%
    \definecolor{cvCream}{RGB}{248,236,224}%
  }{%
  \ifstrequal{#1}{part3}{%
    \definecolor{cvSkyTop}{RGB}{30,20,54}%
    \definecolor{cvSkyMid}{RGB}{92,62,124}%
    \definecolor{cvSkyBot}{RGB}{216,132,140}%
    \definecolor{cvSun}{RGB}{255,214,196}%
    \definecolor{cvFieldTop}{RGB}{128,90,74}%
    \definecolor{cvFieldBot}{RGB}{44,26,30}%
    \definecolor{cvFurrow}{RGB}{34,20,30}%
    \definecolor{cvSil}{RGB}{22,14,22}%
    \definecolor{cvCream}{RGB}{246,238,238}%
  }{%
  \ifstrequal{#1}{part4}{%
    \definecolor{cvSkyTop}{RGB}{12,28,28}%
    \definecolor{cvSkyMid}{RGB}{44,86,72}%
    \definecolor{cvSkyBot}{RGB}{186,168,86}%
    \definecolor{cvSun}{RGB}{240,232,168}%
    \definecolor{cvFieldTop}{RGB}{110,104,50}%
    \definecolor{cvFieldBot}{RGB}{30,34,20}%
    \definecolor{cvFurrow}{RGB}{26,32,20}%
    \definecolor{cvSil}{RGB}{14,20,14}%
    \definecolor{cvCream}{RGB}{240,238,226}%
  }{%
    % 默认: 全书(麦浪金褐)
    \definecolor{cvSkyTop}{RGB}{28,40,32}%
    \definecolor{cvSkyMid}{RGB}{128,110,64}%
    \definecolor{cvSkyBot}{RGB}{232,178,88}%
    \definecolor{cvSun}{RGB}{255,236,180}%
    \definecolor{cvFieldTop}{RGB}{168,124,52}%
    \definecolor{cvFieldBot}{RGB}{62,40,20}%
    \definecolor{cvFurrow}{RGB}{42,26,12}%
    \definecolor{cvSil}{RGB}{28,18,10}%
    \definecolor{cvCream}{RGB}{250,242,226}%
  }}}}}%
}

% --------------------------------------------------------------
%  一键设定:  \CoverSetup{配色}{卷次}{卷名}{副题}{原书页码范围}
% --------------------------------------------------------------
\newcommand{\CoverSetup}[5]{%
  \CoverPalette{#1}%
  \newcommand{\VolNum}{#2}%
  \newcommand{\VolName}{#3}%
  \newcommand{\VolSub}{#4}%
  \newcommand{\VolPages}{#5}%
}

% --------------------------------------------------------------
%  一支麦穗:  \cvWheatEar{x}{底 y}{高}{倾角}
%  茎略弯, 麦粒成对左右斜出, 顶端一芒三针
% --------------------------------------------------------------
\newcommand{\cvWheatEar}[4]{%
  \begin{scope}[shift={(#1,#2)}, rotate=#4]
    \draw[cvSil, line width=0.85mm, line cap=round]
      (0,0) .. controls (0.5,{0.35*#3}) and (0.8,{0.72*#3}) .. ({0.55*#3},#3);
    \foreach \cvK in {0,...,6}{%
      \pgfmathsetmacro{\cvT}{0.50+0.50*\cvK/6}
      \pgfmathsetmacro{\cvGY}{\cvT*#3}
      \pgfmathsetmacro{\cvGX}{0.55*\cvT*\cvT*#3}
      \pgfmathsetmacro{\cvGS}{1.0-0.34*\cvT}
      \begin{scope}[shift={(\cvGX,\cvGY)}, rotate=28]
        \fill[cvSil] (0,{1.9*\cvGS}) ellipse [x radius={0.92*\cvGS}, y radius={2.3*\cvGS}];
      \end{scope}
      \begin{scope}[shift={(\cvGX,\cvGY)}, rotate=-28]
        \fill[cvSil] (0,{-1.9*\cvGS}) ellipse [x radius={0.92*\cvGS}, y radius={2.3*\cvGS}];
      \end{scope}
    }
    \begin{scope}[shift={({0.55*#3},#3)}]
      \fill[cvSil] (0,1.5) ellipse [x radius=0.78, y radius=1.95];
    \end{scope}
    \draw[cvSil, line width=0.26mm, line cap=round]
      ({0.55*#3},{#3+2.3}) -- ({0.55*#3-1.7},{#3+9.5});
    \draw[cvSil, line width=0.26mm, line cap=round]
      ({0.55*#3},{#3+2.3}) -- ({0.55*#3+2.1},{#3+10.5});
    \draw[cvSil, line width=0.26mm, line cap=round]
      ({0.55*#3},{#3+2.3}) -- ({0.55*#3},{#3+11.5});
  \end{scope}
}

% --------------------------------------------------------------
%  封面
% --------------------------------------------------------------
\makeatletter
\renewcommand{\makecover}{%
  \begin{titlepage}
    \thispagestyle{empty}
    \begin{tikzpicture}[remember picture, overlay, x=1mm, y=1mm,
                        shift={(current page.south west)}]
      % ===== 天空 =====
      \shade[top color=cvSkyTop, middle color=cvSkyMid, bottom color=cvSkyBot]
        (0,105) rectangle (210,297);
      % ===== 田野 =====
      \shade[top color=cvFieldTop, bottom color=cvFieldBot]
        (0,0) rectangle (210,105);
      % ===== 落日与光晕 =====
      \begin{scope}
        \clip (0,105) rectangle (210,297);
        \fill[cvSun, opacity=0.14] (152,132) circle (46);
        \fill[cvSun, opacity=0.24] (152,132) circle (30);
      \end{scope}
      \fill[cvSun, opacity=0.95] (152,132) circle (16.5);
      % ===== 犁沟: 向地平线灭点收束 =====
      \begin{scope}
        \clip (0,0) rectangle (210,105);
        \foreach \cvI in {1,...,20}{%
          \pgfmathsetmacro{\cvDX}{3.1*pow(\cvI,1.55)}
          \draw[cvFurrow, opacity=0.30, line width=0.32mm] ({105+\cvDX},0) -- (105,105);
          \draw[cvFurrow, opacity=0.30, line width=0.32mm] ({105-\cvDX},0) -- (105,105);
        }
        \foreach \cvY in {8,18,31,47,68,92}{%
          \draw[cvFurrow, opacity=0.15, line width=0.3mm]
            (0,\cvY) .. controls (60,{\cvY+3.5}) and (150,{\cvY-3}) .. (210,{\cvY+1.5});
        }
      \end{scope}
      % 地平线
      \draw[cvSil, opacity=0.45, line width=0.45mm] (0,105) -- (210,105);
      % ===== 播种者剪影(立于地平线) =====
      \begin{scope}[cvSil]
        \draw[line width=1.9mm, line cap=round, line join=round]
          (55.5,116) -- (50,105)  (55.5,116) -- (60.5,105.5);
        \draw[line width=2.3mm, line cap=round, line join=round]
          (55.5,116) -- (55,123.5);
        \draw[line width=1.6mm, line cap=round, line join=round]
          (55.2,123.2) -- (61.5,126.6) -- (68,129.4);
        \draw[line width=1.5mm, line cap=round, line join=round]
          (55.2,123.2) -- (51.4,119.8) -- (48.6,116.2);
        \fill (55.6,127.3) circle (2.35);
      \end{scope}
      % ===== 种子: 沿贝塞尔弧撒出 =====
      \draw[cvCream, opacity=0.13, line width=0.5mm]
        (68,129.4) .. controls (100,168) and (160,162) .. (196,104);
      % 掌心那一把(尚未撒出)
      \foreach \cvI in {1,...,5}{%
        \pgfmathsetmacro{\cvJX}{1.1*sin(\cvI*53)}
        \pgfmathsetmacro{\cvJY}{1.1*cos(\cvI*83)}
        \fill[cvCream, opacity=0.9] ({68+\cvJX},{129.4+\cvJY}) circle (1.0);
      }
      \foreach \cvI in {1,...,26}{%
        \pgfmathsetmacro{\cvT}{0.035+0.92*\cvI/26}
        \pgfmathsetmacro{\cvU}{1-\cvT}
        \pgfmathsetmacro{\cvSX}{\cvU*\cvU*\cvU*68 + 3*\cvU*\cvU*\cvT*100
                                + 3*\cvU*\cvT*\cvT*160 + \cvT*\cvT*\cvT*196}
        \pgfmathsetmacro{\cvSY}{\cvU*\cvU*\cvU*129.4 + 3*\cvU*\cvU*\cvT*168
                                + 3*\cvU*\cvT*\cvT*162 + \cvT*\cvT*\cvT*104}
        \pgfmathsetmacro{\cvR}{1.55-0.85*\cvT}
        \pgfmathsetmacro{\cvJX}{1.7*sin(\cvI*97)}
        \pgfmathsetmacro{\cvJY}{1.7*cos(\cvI*71)}
        \fill[cvCream, opacity={0.88-0.28*\cvT}] ({\cvSX+\cvJX},{\cvSY+\cvJY}) circle (\cvR);
      }
      % ===== 前景麦穗(两排) =====
      \foreach \cvI in {1,...,11}{%
        \pgfmathsetmacro{\cvWX}{4+19*(\cvI-1)+1.6*sin(\cvI*2.1)}
        \pgfmathsetmacro{\cvWH}{23+7*sin(\cvI*1.7)}
        \pgfmathsetmacro{\cvWL}{3.2*sin(\cvI*2.9)}
        \begin{scope}[opacity=0.45]
          \cvWheatEar{\cvWX}{30}{\cvWH}{\cvWL}
        \end{scope}
      }
      \foreach \cvI in {1,...,10}{%
        \pgfmathsetmacro{\cvWX}{9+20.4*(\cvI-1)+1.5*sin(\cvI*1.3)}
        \pgfmathsetmacro{\cvWH}{41+10*sin(\cvI*2.3)}
        \pgfmathsetmacro{\cvWL}{4.5*sin(\cvI*3.7)}
        \cvWheatEar{\cvWX}{30}{\cvWH}{\cvWL}
      }
      % ===== 压暗一层, 统一色调并保证文字可读 =====
      \fill[black, opacity=0.10] (0,105) rectangle (210,297);
      % ===== 底部信息带 =====
      \fill[cvSil, opacity=0.72] (0,0) rectangle (210,30);
      \draw[cvCream, opacity=0.40, line width=0.4mm] (0,30) -- (210,30);
      % ===== 文字 =====
      \node[anchor=north west, inner sep=0, text=cvCream] at (22,268)
        {\fontsize{38pt}{45pt}\selectfont\headfont \ziju{0.08}\BookTitleCN};
      \node[anchor=north west, inner sep=0, text=cvCream] at (22.6,242)
        {\fontsize{15pt}{20pt}\selectfont\itshape \BookTitleEN};
      \draw[cvCream, opacity=0.55, line width=0.5mm] (22,233) -- (120,233);
      \node[anchor=north west, inner sep=0, text=cvCream] at (22,224)
        {\fontsize{11pt}{15pt}\selectfont 一位数学家对往昔的反思与见证};
      \node[anchor=north west, inner sep=0, text=cvCream, opacity=0.78] at (22.4,208)
        {\fontsize{10pt}{14pt}\selectfont \ziju{0.14}\VolNum};
      \node[anchor=north west, inner sep=0, text=cvCream] at (22,199)
        {\fontsize{30pt}{37pt}\selectfont\headfont \ziju{0.1}\VolName};
      \node[anchor=north west, inner sep=0, text=cvCream, opacity=0.92] at (22,180)
        {\fontsize{13pt}{18pt}\selectfont \VolSub};
      \node[anchor=north west, inner sep=0, text=cvCream] at (22,19)
        {\fontsize{11.5pt}{15pt}\selectfont \OriginalAuthor\quad 著};
      \node[anchor=north west, inner sep=0, text=cvCream, opacity=0.85] at (22,7)
        {\fontsize{9.5pt}{13pt}\selectfont \Translator\ 译\quad ·\quad
         \ifdefempty{\TranslationDate}{\today}{\TranslationDate}};
      \node[anchor=north east, inner sep=0, text=cvCream, opacity=0.85] at (188,13)
        {\fontsize{10pt}{14pt}\selectfont 原书 \VolPages};
    \end{tikzpicture}
  \end{titlepage}%
}
\makeatother
